\documentclass{amsart}
% For dealing with references we use the comment environment
\usepackage{verbatim}
\newenvironment{reference}{\comment}{\endcomment}
%\newenvironment{reference}{}{}
\newenvironment{slogan}{\comment}{\endcomment}
\newenvironment{history}{\comment}{\endcomment}

% For commutative diagrams we use Xy-pic
\usepackage[all]{xy}

% We use 2cell for 2-commutative diagrams.
\xyoption{2cell}
\UseAllTwocells

% We use multicol for the list of chapters between chapters
\usepackage{multicol}

% This is generally recommended for better output
\usepackage{lmodern}
\usepackage[T1]{fontenc}

% For cross-file-references

% Package for hypertext links:
\usepackage{hyperref}

% For any local file, say "hello.tex" you want to link to please

% Theorem environments.
%
\theoremstyle{plain}
\newtheorem{theorem}[subsection]{Theorem}
\newtheorem{proposition}[subsection]{Proposition}
\newtheorem{lemma}[subsection]{Lemma}

\theoremstyle{definition}
\newtheorem{definition}[subsection]{Definition}
\newtheorem{example}[subsection]{Example}
\newtheorem{exercise}[subsection]{Exercise}
\newtheorem{situation}[subsection]{Situation}

\theoremstyle{remark}
\newtheorem{remark}[subsection]{Remark}
\newtheorem{remarks}[subsection]{Remarks}

\numberwithin{equation}{subsection}

% Macros
%
\def\lim{\mathop{\mathrm{lim}}\nolimits}
\def\colim{\mathop{\mathrm{colim}}\nolimits}
\def\Spec{\mathop{\mathrm{Spec}}}
\def\Hom{\mathop{\mathrm{Hom}}\nolimits}
\def\Ext{\mathop{\mathrm{Ext}}\nolimits}
\def\SheafHom{\mathop{\mathcal{H}\!\mathit{om}}\nolimits}
\def\SheafExt{\mathop{\mathcal{E}\!\mathit{xt}}\nolimits}
\def\Sch{\mathit{Sch}}
\def\Mor{\mathop{\mathrm{Mor}}\nolimits}
\def\Ob{\mathop{\mathrm{Ob}}\nolimits}
\def\Sh{\mathop{\mathit{Sh}}\nolimits}
\def\NL{\mathop{N\!L}\nolimits}
\def\CH{\mathop{\mathrm{CH}}\nolimits}
\def\proetale{{pro\text{-}\acute{e}tale}}
\def\etale{{\acute{e}tale}}
\def\QCoh{\mathit{QCoh}}
\def\Ker{\mathop{\mathrm{Ker}}}
\def\Im{\mathop{\mathrm{Im}}}
\def\Coker{\mathop{\mathrm{Coker}}}
\def\Coim{\mathop{\mathrm{Coim}}}

% Boxtimes
%
\DeclareMathSymbol{\boxtimes}{\mathbin}{AMSa}{"02}

%
% Macros for moduli stacks/spaces
%
\def\QCohstack{\mathcal{QC}\!\mathit{oh}}
\def\Cohstack{\mathcal{C}\!\mathit{oh}}
\def\Spacesstack{\mathcal{S}\!\mathit{paces}}
\def\Quotfunctor{\mathrm{Quot}}
\def\Hilbfunctor{\mathrm{Hilb}}
\def\Curvesstack{\mathcal{C}\!\mathit{urves}}
\def\Polarizedstack{\mathcal{P}\!\mathit{olarized}}
\def\Complexesstack{\mathcal{C}\!\mathit{omplexes}}
% \Pic is the operator that assigns to X its picard group, usage \Pic(X)
% \Picardstack_{X/B} denotes the Picard stack of X over B
% \Picardfunctor_{X/B} denotes the Picard functor of X over B
\def\Pic{\mathop{\mathrm{Pic}}\nolimits}
\def\Picardstack{\mathcal{P}\!\mathit{ic}}
\def\Picardfunctor{\mathrm{Pic}}
\def\Deformationcategory{\mathcal{D}\!\mathit{ef}}

\setlength{\emergencystretch}{3em}
\begin{document}
\title[Stacks Project, Tag 07SU]{1743 --- Every quasi-compact quasi-separated algebraic space over the integers is an affine-transition inverse limit of finite-type quasi-separated integer algebraic spaces}
\maketitle
\noindent Copyright (C) 2005--2025 Johan de Jong. Stacks Project authors. Source: \href{https://stacks.math.columbia.edu/tag/07SU}{Tag 07SU}.

\noindent Editorial difficulty note (not author text). Difficulty H3: The theorem builds finite-type integer models of an arbitrary qcqs algebraic space through finite Nisnevich filtration, relative scheme approximation and explicit gluing of finite-stage etale equivalence relations. The new relation is manufactured by adjoining the complement diagonal to a descended overlap relation, then cartesian transitions and quotient limits are recovered. This non-schematic approximation/effectiveness construction is substantially beyond routine scheme descent or formal parameter changes..

\noindent Editorial provenance note (not author text). History: excerpt from revision \texttt{a0444\allowbreak 6e57e\allowbreak c1fbc\allowbreak 252a8\allowbreak 71afc\allowbreak ec775\allowbreak 2fb28\allowbreak 07b14}, spaces-limits.tex, lines 2110--2260. Reference presentation was adapted; original-author context and core support blocks were retained or added. The mathematical wording is unchanged. Licensed under GNU FDL 1.2 or later; no invariant sections or cover texts. See ../licenses/COPYING. Context blocks reproduce author definitions and supporting results; they are not additional counted problems.

\noindent
Let $S$ be a scheme.
Let $(U, R, s, t, c)$ be a groupoid scheme over $S$.
Assume $U = \Spec(A)$, and $R = \Spec(B)$ are affine.
In this case we get two ring maps
$s^\sharp, t^\sharp : A \longrightarrow B$.
Let $C$ be the equalizer of $s^\sharp$ and $t^\sharp$. In a formula
\begin{equation}
\label{groupoids-equation-invariants}
C = \{a \in A \mid t^\sharp(a) = s^\sharp(a) \}.
\end{equation}
We will sometimes call this the {\it ring of $R$-invariant functions on $U$}.
What properties does $M = \Spec(C)$ have? The first observation is
that the diagram
$$
\xymatrix{
R \ar[r]_s \ar[d]_t & U \ar[d] \\
U \ar[r] & M
}
$$
is commutative, i.e., the morphism $U \to M$ equalizes $s, t$.
Moreover, if $T$ is any affine scheme, and if $U \to T$ is
a morphism which equalizes $s, t$, then $U \to T$ factors through $U \to M$.
In other words, $U \to M$ is a coequalizer in the category of affine schemes.

\begin{definition}
\label{groupoids-definition-quotient-sheaf}
Let $\tau$, $S$, and the pre-relation $j : R \to U \times_S U$ be as above.
In this setting the {\it quotient sheaf $U/R$} associated
to $j$ is the sheafification of the presheaf
(\ref{groupoids-equation-quotient-presheaf}) in the $\tau$-topology.
If $j : R \to U \times_S U$ comes from the action of a group scheme
$G/S$ on $U$ as in Lemma \href{https://stacks.math.columbia.edu/tag/0234}{lemma groupoid from action (Tag 0234)} then we
sometimes denote the quotient sheaf $U/G$.
\end{definition}

\begin{definition}
\label{groupoids-definition-groupoid}
Let $S$ be a scheme.
\begin{enumerate}
\item A {\it groupoid scheme over $S$}, or simply a
{\it groupoid over $S$} is a
quintuple $(U, R, s, t, c)$ where
$U$ and $R$ are schemes over $S$, and
$s, t : R \to U$ and $c : R \times_{s, U, t} R \to R$
are morphisms of schemes over $S$ with the
following property: For any scheme
$T$ over $S$ the quintuple
$$
(U(T), R(T), s, t, c)
$$
is a groupoid category in the sense described above.
\item A {\it morphism
$f : (U, R, s, t, c) \to (U', R', s', t', c')$
of groupoid schemes over $S$} is given by morphisms
of schemes $f : U \to U'$ and $f : R \to R'$ with the
following property:  For any scheme
$T$ over $S$ the maps $f$ define a functor from the
groupoid category $(U(T), R(T), s, t, c)$ to the
groupoid category $(U'(T), R'(T), s', t', c')$.
\end{enumerate}
\end{definition}

\begin{definition}
\label{groupoids-definition-equivalence-relation}
Let $S$ be a scheme. Let $U$ be a scheme over $S$.
\begin{enumerate}
\item A {\it pre-relation} on $U$ over $S$ is any morphism
of schemes $j : R \to U \times_S U$. In this case we set
$t = \text{pr}_0 \circ j$ and $s = \text{pr}_1 \circ j$, so
that $j = (t, s)$.
\item A {\it relation} on $U$ over $S$ is a monomorphism
of schemes $j : R \to U \times_S U$.
\item A {\it pre-equivalence relation} is a pre-relation
$j : R \to U \times_S U$ such that the image of
$j : R(T) \to U(T) \times U(T)$ is an equivalence relation for
all $T/S$.
\item We say a morphism $R \to U \times_S U$ of schemes is
an {\it equivalence relation on $U$ over $S$}
if and only if for every scheme $T$ over $S$ the $T$-valued
points of $R$ define an equivalence relation
on the set of $T$-valued points of $U$.
\end{enumerate}
\end{definition}

\begin{definition}
\label{categories-definition-groupoid}
A {\it groupoid} is a category where every morphism is an isomorphism.
\end{definition}

\begin{situation}
\label{limits-situation-descent}
Let $S = \lim_{i \in I} S_i$ be the limit of a directed system of schemes
with affine transition morphisms $f_{i'i} : S_{i'} \to S_i$
(Lemma \ref{limits-lemma-directed-inverse-system-has-limit}).
We assume that $S_i$ is quasi-compact and quasi-separated for all $i \in I$.
We denote $f_i : S \to S_i$ the projection.
We also choose an element $0 \in I$.
\end{situation}

\begin{definition}
\label{properties-definition-ample}
\begin{reference}
\href{https://stacks.math.columbia.edu/bibliography}{Bibliography: EGA (II Definition 4.5.3)}
\end{reference}
Let $X$ be a scheme.
Let $\mathcal{L}$ be an invertible $\mathcal{O}_X$-module.
We say $\mathcal{L}$ is {\it ample} if
\begin{enumerate}
\item $X$ is quasi-compact, and
\item for every $x \in X$ there exists an $n \geq 1$
and $s \in \Gamma(X, \mathcal{L}^{\otimes n})$ such
that $x \in X_s$ and $X_s$ is affine.
\end{enumerate}
\end{definition}

\begin{definition}
\label{schemes-definition-reduced-induced-scheme}
Let $X$ be a scheme. Let $Z \subset X$ be a closed subset.
A {\it scheme structure on $Z$} is given by a closed subscheme $Z'$ of
$X$ whose underlying set is equal to $Z$. We often say
``let $(Z, \mathcal{O}_Z)$ be a scheme structure on $Z$'' to
indicate this. The {\it reduced induced scheme structure}
on $Z$ is the one constructed in Lemma \href{https://stacks.math.columbia.edu/tag/01J3}{lemma reduced closed subscheme (Tag 01J3)}.
The {\it reduction $X_{red}$ of $X$} is the reduced induced scheme
structure on $X$ itself.
\end{definition}

\begin{situation}
\label{situation-descent}
Let $S$ be a scheme. Let $X = \lim_{i \in I} X_i$ be the limit of a directed
inverse system of algebraic spaces over $S$ with affine transition morphisms
(Lemma \ref{lemma-directed-inverse-system-has-limit}).
We assume that $X_i$ is quasi-compact and quasi-separated for all $i \in I$.
We also choose an element $0 \in I$.
\end{situation}

\begin{definition}
\label{decent-spaces-definition-very-reasonable}
Let $S$ be a scheme.
Let $X$ be an algebraic space over $S$.
\begin{enumerate}
\item We say $X$ is {\it decent} if for every point $x \in X$ the equivalent
conditions of
Lemma \href{https://stacks.math.columbia.edu/tag/03JV}{lemma UR finite above x (Tag 03JV)}
hold, in other words property $(\gamma)$ of
Lemma \ref{decent-spaces-lemma-bounded-fibres}
holds.
\item We say $X$ is {\it reasonable} if the equivalent conditions of
Lemma \href{https://stacks.math.columbia.edu/tag/03JT}{lemma U universally bounded (Tag 03JT)}
hold, in other words property $(\delta)$ of
Lemma \ref{decent-spaces-lemma-bounded-fibres}
holds.
\item We say $X$ is {\it very reasonable} if the equivalent conditions of
Lemma \href{https://stacks.math.columbia.edu/tag/03IH}{lemma characterize very reasonable (Tag 03IH)}
hold, i.e., property $(\epsilon)$ of
Lemma \ref{decent-spaces-lemma-bounded-fibres}
holds.
\end{enumerate}
\end{definition}

\begin{definition}
\label{spaces-definition-presentation}
Let $F$ be an algebraic space over $S$.
A {\it presentation} of $F$ is given by a scheme
$U$ over $S$ and an \'etale equivalence relation $R$ on $U$ over $S$, and
a surjective \'etale morphism $U \to F$ such that $R = U \times_F U$.
\end{definition}

\begin{definition}
\label{groupoids-definition-restrict-relation}
Let $S$ be a scheme.
Let $U$ be a scheme over $S$.
Let $j : R \to U \times_S U$ be a pre-relation.
Let $g : U' \to U$ be a morphism of schemes.
The pre-relation $j' : R' \to U' \times_S U'$ is called
the {\it restriction}, or {\it pullback} of the pre-relation $j$ to $U'$.
In this situation we sometimes write $R' = R|_{U'}$.
\end{definition}

\begin{definition}
\label{spaces-morphisms-definition-locally-finite-presentation}
Let $S$ be a scheme.
Let $f : X \to Y$ be a morphism of algebraic spaces over $S$.
\begin{enumerate}
\item We say $f$ is {\it locally of finite presentation} if
the equivalent conditions of
Lemma \href{https://stacks.math.columbia.edu/tag/03MJ}{lemma local source target (Tag 03MJ)}
hold with $\mathcal{P} =$``locally of finite presentation''.
\item Let $x \in |X|$. We say $f$ is of {\it finite presentation at $x$}
if there exists an open neighbourhood $X' \subset X$ of $x$ such
that $f|_{X'} : X' \to Y$ is locally of finite
presentation\footnote{It seems awkward to use ``locally of finite presentation
at $x$'', but the current terminology may be misleading in the sense that
``of finite presentation at $x$'' does {\bf not} mean that there is
an open neighbourhood $X' \subset X$ such that $f|_{X'}$ is of finite
presentation.}.
\item A morphism of algebraic spaces $f : X \to Y$ is
{\it of finite presentation}
if it is locally of finite presentation, quasi-compact and
quasi-separated.
\end{enumerate}
\end{definition}

\begin{definition}
\label{definition-locally-finite-presentation}
Let $S$ be a scheme.
\begin{enumerate}
\item A functor $F : (\Sch/S)_{fppf}^{opp} \to \textit{Sets}$
is said to be {\it limit preserving} or {\it locally of finite presentation} if
for every affine scheme $T$ over $S$ which is a limit $T = \lim T_i$
of a directed inverse system of affine schemes $T_i$ over $S$, we have
$$
F(T) = \colim F(T_i).
$$
We sometimes say that $F$ is {\it locally of finite presentation over $S$}.
\item Let $F, G : (\Sch/S)_{fppf}^{opp} \to \textit{Sets}$.
A transformation of functors $a : F \to G$
is {\it limit preserving} or {\it locally of finite presentation}
if for every scheme $T$ over $S$ and every $y \in G(T)$ the functor
$$
F_y : (\Sch/T)_{fppf}^{opp} \longrightarrow \textit{Sets}, \quad
T'/T \longmapsto \{x \in F(T') \mid a(x) = y|_{T'}\}
$$
is locally of finite presentation over $T$\footnote{The characterization (2) in
Lemma \href{https://stacks.math.columbia.edu/tag/06BC}{lemma characterize relative limit preserving (Tag 06BC)}
may be easier to parse.}. We sometimes say that
$F$ is {\it relatively limit preserving} over $G$.
\end{enumerate}
\end{definition}

\begin{situation}
\label{limits-situation-descent-property}
Let $S = \lim S_i$ be a limit of a directed system of schemes
with affine transition morphisms
(Lemma \ref{limits-lemma-directed-inverse-system-has-limit}).
Let $0 \in I$ and let $f_0 : X_0 \to Y_0$ be a morphism of schemes over $S_0$.
Assume $S_0$, $X_0$, $Y_0$ are quasi-compact and quasi-separated.
Let $f_i : X_i \to Y_i$ be the base change of $f_0$ to $S_i$ and
let $f : X \to Y$ be the base change of $f_0$ to $S$.
\end{situation}

\begin{situation}
\label{situation-descent-property}
Let $S$ be a scheme. Let $B = \lim B_i$ be a limit of a directed inverse system
of algebraic spaces over $S$ with affine transition morphisms
(Lemma \ref{lemma-directed-inverse-system-has-limit}).
Let $0 \in I$ and let $f_0 : X_0 \to Y_0$ be a morphism of algebraic spaces
over $B_0$. Assume $B_0$, $X_0$, $Y_0$ are quasi-compact and quasi-separated.
Let $f_i : X_i \to Y_i$ be the base change of $f_0$ to $B_i$ and
let $f : X \to Y$ be the base change of $f_0$ to $B$.
\end{situation}

\begin{lemma}
\label{lemma-descend-finite-presentation}
Let $S$ be a scheme. Let $I$ be a directed set.
Let $(X_i, f_{ii'})$ be an inverse system over $I$ of algebraic spaces
over $S$. Assume
\begin{enumerate}
\item the morphisms $f_{ii'} : X_i \to X_{i'}$ are affine,
\item the spaces $X_i$ are quasi-compact and quasi-separated.
\end{enumerate}
Let $X = \lim_i X_i$. Then the category of algebraic spaces
of finite presentation over $X$ is the colimit over $I$ of the
categories of algebraic spaces of finite presentation over $X_i$.
\end{lemma}

\begin{proof}
Pick $0 \in I$. Choose a surjective \'etale morphism $U_0 \to X_0$ where
$U_0$ is an affine scheme (Properties of Spaces, Lemma
\href{https://stacks.math.columbia.edu/tag/03H6}{lemma quasi compact affine cover (Tag 03H6)}).
Set $U_i = X_i \times_{X_0} U_0$. Set $R_0 = U_0 \times_{X_0} U_0$ and
$R_i = R_0 \times_{X_0} X_i$. Denote $s_i, t_i : R_i \to U_i$ and
$s, t : R \to U$ the two projections. In the proof of
Lemma \ref{lemma-directed-inverse-system-has-limit} we have
seen that there exists a presentation $X = U/R$ with
$U = \lim U_i$ and $R = \lim R_i$. Note that $U_i$ and $U$ are affine and
that $R_i$ and $R$ are quasi-compact and separated (as $X_i$ is
quasi-separated). Let $Y$ be an algebraic space over $S$ and let
$Y \to X$ be a morphism of finite presentation. Set $V = U \times_X Y$.
This is an algebraic space of finite presentation over $U$.
Choose an affine scheme $W$ and a surjective \'etale morphism $W \to V$.
Then $W \to Y$ is surjective \'etale as well. Set $R' = W \times_Y W$
so that $Y = W/R'$ (see Spaces, Section \href{https://stacks.math.columbia.edu/tag/0261}{section presentations (Tag 0261)}).
Note that $W$ is a scheme of finite presentation over $U$ and that $R'$
is a scheme of finite presentation over $R$ (details omitted).
By Limits, Lemma \ref{limits-lemma-descend-finite-presentation}
we can find an index $i$ and a morphism of schemes $W_i \to U_i$ of
finite presentation whose base change to $U$ gives $W \to U$. Similarly
we can find, after possibly increasing $i$, a scheme $R'_i$ of finite
presentation over $R_i$ whose base change to $R$ is $R'$.
The projection morphisms $s', t' : R' \to W$ are morphisms over
the projection morphisms $s, t : R \to U$. Hence we can view $s'$,
resp.\ $t'$ as a morphism between schemes of finite presentation over
$U$ (with structure morphism $R' \to U$ given by $R' \to R$ followed
by $s$, resp.\ $t$). Hence we can apply
Limits, Lemma \ref{limits-lemma-descend-finite-presentation}
again to see that, after possibly increasing $i$, there exist
morphisms $s'_i, t'_i : R'_i \to W_i$, whose base change to $U$
is $S', t'$. By Limits, Lemmas \ref{limits-lemma-descend-etale} and
\ref{limits-lemma-descend-monomorphism}
we may assume that $s'_i, t'_i$ are \'etale and that
$j'_i : R'_i \to W_i \times_{X_i} W_i$ is a monomorphism (here we
view $j'_i$ as a morphism of schemes of finite presentation over $U_i$ via
one of the projections -- it doesn't matter which one). Setting
$Y_i = W_i/R'_i$ (see Spaces, Theorem \ref{spaces-theorem-presentation})
we obtain an algebraic space of finite presentation
over $X_i$ whose base change to $X$ is isomorphic to $Y$.

\medskip\noindent
This shows that every algebraic space of finite presentation over $X$ comes
from an algebraic space of finite presentation over some $X_i$, i.e.,
it shows that the functor of the lemma is essentially surjective. To
show that it is fully faithful, consider an index $0 \in I$ and two
algebraic spaces $Y_0, Z_0$ of finite presentation over $X_0$.
Set $Y_i = X_i \times_{X_0} Y_0$, $Y = X \times_{X_0} Y_0$,
$Z_i = X_i \times_{X_0} Z_0$, and $Z = X \times_{X_0} Z_0$. Let
$\alpha : Y \to Z$ be a morphism of algebraic spaces over $X$.
Choose a surjective \'etale morphism $V_0 \to Y_0$ where $V_0$ is
an affine scheme. Set $V_i = V_0 \times_{Y_0} Y_i$ and
$V = V_0 \times_{Y_0} Y$ which are affine schemes endowed with
surjective \'etale morphisms to $Y_i$ and $Y$. The composition
$V \to Y \to Z \to Z_0$ comes from a (essentially unique) morphism
$V_i \to Z_0$ for some $i \geq 0$ by
Proposition \ref{proposition-characterize-locally-finite-presentation}
(applied to $Z_0 \to X_0$ which is of finite presentation by assumption).
After increasing $i$ the two compositions
$$
V_i \times_{Y_i} V_i \to V_i \to Z_0
$$
are equal as this is true in the limit. Hence we obtain a (essentially unique)
morphism $Y_i \to Z_0$. Since this is a morphism over $X_0$
it induces a morphism into $Z_i = Z_0 \times_{X_0} X_i$ as desired.
\end{proof}


\begin{lemma}
\label{lemma-descend-etale}
With notation and assumptions as in
Situation \ref{situation-descent-property}. If
\begin{enumerate}
\item $f$ is \'etale,
\item $f_0$ is locally of finite presentation,
\end{enumerate}
then $f_i$ is \'etale for some $i \geq 0$.
\end{lemma}

\begin{proof}
Choose an affine scheme $V_0$ and a surjective \'etale morphism
$V_0 \to Y_0$. Choose an affine scheme $U_0$ and a surjective \'etale
morphism $U_0 \to V_0 \times_{Y_0} X_0$. Diagram
$$
\xymatrix{
U_0 \ar[d] \ar[r] & V_0 \ar[d] \\
X_0 \ar[r] & Y_0
}
$$
The vertical arrows are surjective and \'etale by construction.
We can base change this diagram to $B_i$ or $B$ to get
$$
\vcenter{
\xymatrix{
U_i \ar[d] \ar[r] & V_i \ar[d] \\
X_i \ar[r] & Y_i
}
}
\quad\text{and}\quad
\vcenter{
\xymatrix{
U \ar[d] \ar[r] & V \ar[d] \\
X \ar[r] & Y
}
}
$$
Note that $U_i, V_i, U, V$ are affine schemes,
the vertical morphisms are surjective \'etale, and the limit of the
morphisms $U_i \to V_i$ is $U \to V$. Recall that $X_i \to Y_i$ is \'etale
if and only if $U_i \to V_i$ is
\'etale and similarly $X \to Y$ is \'etale if and only if
$U \to V$ is \'etale
(Morphisms of Spaces, Lemma \href{https://stacks.math.columbia.edu/tag/03XT}{lemma etale local (Tag 03XT)}).
Since $f_0$ is locally of finite
presentation, so is the morphism $U_0 \to V_0$. Hence the lemma follows
from Limits, Lemma \ref{limits-lemma-descend-etale}.
\end{proof}


\begin{lemma}
\label{lemma-directed-inverse-system-has-limit}
Let $S$ be a scheme. Let $I$ be a directed set.
Let $(X_i, f_{ii'})$ be an inverse system over $I$
in the category of algebraic spaces over $S$.
If the morphisms $f_{ii'} : X_i \to X_{i'}$ are affine, then the
limit $X = \lim_i X_i$ (as an fppf sheaf) is an algebraic space.
Moreover,
\begin{enumerate}
\item each of the morphisms $f_i : X \to X_i$ is affine,
\item for any $i \in I$ and any morphism of algebraic spaces
$T \to X_i$ we have
$$
X \times_{X_i} T = \lim_{i' \geq i} X_{i'} \times_{X_i} T.
$$
as algebraic spaces over $S$.
\end{enumerate}
\end{lemma}

\begin{proof}
Part (2) is a formal consequence of the existence of the
limit $X = \lim X_i$ as an algebraic space over $S$.
Choose an element $0 \in I$ (this is possible as a directed set is nonempty).
Choose a scheme $U_0$ and a surjective
\'etale morphism $U_0 \to X_0$. Set $R_0 = U_0 \times_{X_0} U_0$
so that $X_0 = U_0/R_0$. For $i \geq 0$ set
$U_i = X_i \times_{X_0} U_0$ and
$R_i = X_i \times_{X_0} R_0 = U_i \times_{X_i} U_i$.
By Limits, Lemma \ref{limits-lemma-directed-inverse-system-has-limit}
we see that $U = \lim_{i \geq 0} U_i$ and $R = \lim_{i \geq 0} R_i$
are schemes. Moreover, the two morphisms $s, t : R \to U$ are the base
change of the two projections $R_0 \to U_0$ by the morphism
$U \to U_0$, in particular \'etale. The morphism $R \to U \times_S U$
defines an equivalence relation as directed a limit of equivalence relations
is an equivalence relation. Hence the morphism
$R \to U \times_S U$ is an \'etale equivalence relation. We claim that
the natural map
\begin{equation}
\label{equation-isomorphism-sheaves}
U/R \longrightarrow \lim X_i
\end{equation}
is an isomorphism of fppf sheaves on the category of schemes over $S$.
The claim implies $X = \lim X_i$ is an algebraic
space by Spaces, Theorem \ref{spaces-theorem-presentation}.

\medskip\noindent
Let $Z$ be a scheme and let $a : Z \to \lim X_i$ be a morphism.
Then $a = (a_i)$ where $a_i : Z \to X_i$. Set $W_0 = Z \times_{a_0, X_0} U_0$.
Note that $W_0 = Z \times_{a_i, X_i} U_i$ for all $i \geq 0$ by our
choice of $U_i \to X_i$ above. Hence we obtain a morphism
$W_0 \to \lim_{i \geq 0} U_i = U$. Since $W_0 \to Z$ is surjective
and \'etale, we conclude that (\ref{equation-isomorphism-sheaves})
is a surjective map of sheaves. Finally, suppose that
$Z$ is a scheme and that $a, b : Z \to U/R$ are two morphisms
which are equalized by (\ref{equation-isomorphism-sheaves}).
We have to show that $a = b$.
After replacing $Z$ by the members of an fppf covering
we may assume there exist morphisms $a', b' : Z \to U$ which
give rise to $a$ and $b$. The condition that $a, b$ are
equalized by (\ref{equation-isomorphism-sheaves}) means that
for each $i \geq 0$ the compositions $a_i', b_i' : Z \to U \to U_i$
are equal as morphisms into $U_i/R_i = X_i$. Hence
$(a_i', b_i') : Z \to U_i \times_S U_i$ factors through
$R_i$, say by some morphism $c_i : Z \to R_i$. Since
$R = \lim_{i \geq 0} R_i$ we see that $c = \lim c_i : Z \to R$
is a morphism which shows that $a, b$ are equal as morphisms
of $Z$ into $U/R$.

\medskip\noindent
Part (1) follows as we have seen above that
$U_i \times_{X_i} X = U$ and $U \to U_i$ is affine by
construction.
\end{proof}


\begin{lemma}
\label{lemma-limit-is-scheme}
Notation and assumptions as in Situation \ref{situation-descent}.
If $X$ is a scheme, then there exists an $i$ such that $X_i$ is a scheme.
\end{lemma}

\begin{proof}
Choose a finite affine open covering $X = \bigcup W_j$.
By Lemma \ref{lemma-descend-opens}
we can find an $i \in I$ and open subspaces $W_{j, i} \subset X_i$
whose base change to $X$ is $W_j \to X$. By
Lemma \ref{lemma-limit-is-affine} we may assume that
each $W_{j, i}$ is an affine scheme. This means that $X_i$
is a scheme (see for example
Properties of Spaces, Section \href{https://stacks.math.columbia.edu/tag/03JG}{section schematic (Tag 03JG)}).
\end{proof}


\begin{lemma}
\label{lemma-limit-is-affine}
Notation and assumptions as in Situation \ref{situation-descent}.
If $X$ is affine, then there exists an $i$ such that $X_i$ is affine.
\end{lemma}

\begin{proof}
Choose $0 \in I$. Choose an affine scheme $U_0$ and a surjective
\'etale morphism $U_0 \to X_0$. Set $U = U_0 \times_{X_0} X$
and $U_i = U_0 \times_{X_0} X_i$ for $i \geq 0$. Since the transition
morphisms are affine, the algebraic spaces $U_i$ and $U$ are affine.
Thus $U \to X$ is an \'etale morphism of affine schemes. Hence we
can write $X = \Spec(A)$, $U = \Spec(B)$ and
$$
B = A[x_1, \ldots, x_n]/(g_1, \ldots, g_n)
$$
such that $\Delta = \det(\partial g_\lambda/\partial x_\mu)$ is invertible
in $B$, see Algebra, Lemma \href{https://stacks.math.columbia.edu/tag/00U9}{lemma etale standard smooth (Tag 00U9)}.
Set $A_i = \mathcal{O}_{X_i}(X_i)$. We have $A = \colim A_i$ by
Lemma \ref{lemma-descend-section}. After increasing $0$ we may assume
we have $g_{1, i}, \ldots, g_{n, i} \in A_i[x_1, \ldots, x_n]$ mapping to
$g_1, \ldots, g_n$. Set
$$
B_i = A_i[x_1, \ldots, x_n]/(g_{1, i}, \ldots, g_{n, i})
$$
for all $i \geq 0$. Increasing $0$ if necessary we may assume that
$\Delta_i = \det(\partial g_{\lambda, i}/\partial x_\mu)$ is invertible
in $B_i$ for all $i \geq 0$. Thus $A_i \to B_i$ is an \'etale ring map.
After increasing $0$ we may assume also that
$\Spec(B_i) \to \Spec(A_i)$ is surjective, see
Limits, Lemma \href{https://stacks.math.columbia.edu/tag/07RR}{lemma descend surjective (Tag 07RR)}. Increasing
$0$ yet again we may choose elements
$h_{1, i}, \ldots, h_{n, i} \in \mathcal{O}_{U_i}(U_i)$ which map to the
classes of $x_1, \ldots, x_n$ in $B = \mathcal{O}_U(U)$ and such that
$g_{\lambda, i}(h_{\nu, i}) = 0$ in $\mathcal{O}_{U_i}(U_i)$. Thus
we obtain a commutative diagram
\begin{equation}
\label{equation-to-show-cartesian}
\vcenter{
\xymatrix{
X_i \ar[d] & U_i \ar[l] \ar[d] \\
\Spec(A_i) & \Spec(B_i) \ar[l]
}
}
\end{equation}
By construction $B_i = B_0 \otimes_{A_0} A_i$ and
$B = B_0 \otimes_{A_0} A$. Consider the morphism
$$
f_0 : U_0 \longrightarrow X_0 \times_{\Spec(A_0)} \Spec(B_0)
$$
This is a morphism of quasi-compact and quasi-separated algebraic spaces
representable, separated and \'etale over $X_0$. The base change of $f_0$
to $X$ is an isomorphism by our choices. Hence
Lemma \ref{lemma-descend-equality}
guarantees that there exists an $i$ such that the base change of $f_0$
to $X_i$ is an isomorphism, in other words the diagram
(\ref{equation-to-show-cartesian}) is cartesian. Thus
Descent, Lemma \href{https://stacks.math.columbia.edu/tag/02W5}{lemma descent data sheaves (Tag 02W5)}
applied to the fppf covering $\{\Spec(B_i) \to \Spec(A_i)\}$
combined with Descent, Lemma \href{https://stacks.math.columbia.edu/tag/0245}{lemma affine (Tag 0245)}
give that $X_i \to \Spec(A_i)$ is representable by a scheme
affine over $\Spec(A_i)$ as desired. (Of course it then also follows
that $X_i = \Spec(A_i)$ but we don't need this.)
\end{proof}


\begin{lemma}
\label{lemma-descend-section}
Notation and assumptions as in Situation \ref{situation-descent}.
Suppose that $\mathcal{F}_0$ is a quasi-coherent sheaf on $X_0$.
Set $\mathcal{F}_i = f_{0i}^*\mathcal{F}_0$ for $i \geq 0$ and set
$\mathcal{F} = f_0^*\mathcal{F}_0$. Then
$$
\Gamma(X, \mathcal{F}) = \colim_{i \geq 0} \Gamma(X_i, \mathcal{F}_i)
$$
\end{lemma}

\begin{proof}
Choose a surjective \'etale morphism $U_0 \to X_0$ where $U_0$ is an affine
scheme (Properties of Spaces, Lemma
\href{https://stacks.math.columbia.edu/tag/03H6}{lemma quasi compact affine cover (Tag 03H6)}).
Set $U_i = X_i \times_{X_0} U_0$.
Set $R_0 = U_0 \times_{X_0} U_0$ and $R_i = R_0 \times_{X_0} X_i$.
In the proof of Lemma \ref{lemma-directed-inverse-system-has-limit} we have
seen that there exists a presentation $X = U/R$ with
$U = \lim U_i$ and $R = \lim R_i$.
Note that $U_i$ and $U$ are affine and that $R_i$ and $R$ are
quasi-compact and separated (as $X_i$ is quasi-separated). Hence
Limits, Lemma \ref{limits-lemma-descend-section}
implies that
$$
\mathcal{F}(U) = \colim \mathcal{F}_i(U_i)
\quad\text{and}\quad
\mathcal{F}(R) = \colim \mathcal{F}_i(R_i).
$$
The lemma follows as
$\Gamma(X, \mathcal{F}) = \Ker(\mathcal{F}(U) \to \mathcal{F}(R))$
and similarly
$\Gamma(X_i, \mathcal{F}_i) =
\Ker(\mathcal{F}_i(U_i) \to \mathcal{F}_i(R_i))$
\end{proof}


\begin{lemma}
\label{lemma-descend-opens}
Notation and assumptions as in Situation \ref{situation-descent}.
For any quasi-compact open subspace $U \subset X$ there exists an $i$
and a quasi-compact open $U_i \subset X_i$ whose inverse image in $X$ is $U$.
\end{lemma}

\begin{proof}
Follows formally from the construction of limits in
Lemma \ref{lemma-directed-inverse-system-has-limit}
and the corresponding result for schemes:
Limits, Lemma \ref{limits-lemma-descend-opens}.
\end{proof}


\begin{proposition}
\label{proposition-characterize-locally-finite-presentation}
Let $S$ be a scheme. Let $f : X \to Y$ be a morphism of algebraic
spaces over $S$. The following are equivalent:
\begin{enumerate}
\item The morphism $f$ is a morphism of algebraic spaces which is
locally of finite presentation, see
Morphisms of Spaces,
Definition \ref{spaces-morphisms-definition-locally-finite-presentation}.
\item The morphism $f : X \to Y$ is limit preserving as
a transformation of functors, see
Definition \ref{definition-locally-finite-presentation}.
\end{enumerate}
\end{proposition}

\begin{proof}
Assume (1). Let $T$ be a scheme and let $y \in Y(T)$. We have to show that
$T \times_Y X$ is limit preserving over $T$ in the sense of
Definition \ref{definition-locally-finite-presentation}.
Hence we are reduced to proving that if $X$ is an algebraic space which
is locally of finite presentation over $S$ as an algebraic space, then it
is limit preserving as a functor
$X : (\Sch/S)_{fppf}^{opp} \to \textit{Sets}$.
To see this choose a presentation $X = U/R$, see
Spaces, Definition \ref{spaces-definition-presentation}.
It follows from
Morphisms of Spaces,
Definition \ref{spaces-morphisms-definition-locally-finite-presentation}
that both $U$ and $R$ are schemes which are locally of finite presentation
over $S$. Hence by
Limits, Proposition
\href{https://stacks.math.columbia.edu/tag/01ZC}{proposition characterize locally finite presentation (Tag 01ZC)}
we have
$$
U(T) = \colim U(T_i), \quad
R(T) = \colim R(T_i)
$$
whenever $T = \lim_i T_i$ in $(\Sch/S)_{fppf}$. It follows
that the presheaf
$$
(\Sch/S)_{fppf}^{opp} \longrightarrow \textit{Sets}, \quad
W \longmapsto U(W)/R(W)
$$
is limit preserving. Hence by
Lemma \ref{lemma-sheafify-finite-presentation}
its sheafification $X = U/R$ is limit preserving too.

\medskip\noindent
Assume (2). Choose a scheme $V$ and a surjective \'etale morphism
$V \to Y$. Next, choose a scheme $U$ and a surjective \'etale morphism
$U \to V \times_Y X$. By
Lemma \href{https://stacks.math.columbia.edu/tag/049M}{lemma base change locally finite presentation (Tag 049M)}
the transformation of functors $V \times_Y X \to V$ is limit preserving. By
Morphisms of Spaces,
Lemma \href{https://stacks.math.columbia.edu/tag/0468}{lemma etale locally finite presentation (Tag 0468)}
the morphism of algebraic spaces $U \to V \times_Y X$ is locally
of finite presentation, hence limit preserving as
a transformation of functors by the first part of the proof. By
Lemma \href{https://stacks.math.columbia.edu/tag/049L}{lemma composition locally finite presentation (Tag 049L)}
the composition $U \to V \times_Y X \to V$ is limit preserving
as a transformation of functors. Hence
the morphism of schemes $U \to V$ is locally of finite presentation by
Limits, Proposition
\href{https://stacks.math.columbia.edu/tag/01ZC}{proposition characterize locally finite presentation (Tag 01ZC)}
(modulo a set theoretic remark, see last paragraph of the proof).
This means, by definition, that (1) holds.

\medskip\noindent
Set theoretic remark. Let $U \to V$ be a morphism of
$(\Sch/S)_{fppf}$. In the statement of
Limits, Proposition
\href{https://stacks.math.columbia.edu/tag/01ZC}{proposition characterize locally finite presentation (Tag 01ZC)}
we characterize $U \to V$ as being locally of finite presentation
if for {\it all} directed inverse systems $(T_i, f_{ii'})$ of affine schemes
over $V$ we have $U(T) = \colim V(T_i)$, but in the current setting
we may only consider affine schemes $T_i$ over $V$ which are (isomorphic to)
an object of $(\Sch/S)_{fppf}$. So we have to make sure that there
are enough affines in $(\Sch/S)_{fppf}$ to make the proof work.
Inspecting the proof of (2) $\Rightarrow$ (1) of
Limits, Proposition
\href{https://stacks.math.columbia.edu/tag/01ZC}{proposition characterize locally finite presentation (Tag 01ZC)}
we see that the question reduces to the case that $U$ and $V$ are affine.
Say $U = \Spec(A)$ and $V = \Spec(B)$. By construction
of $(\Sch/S)_{fppf}$ the spectrum of any ring of cardinality
$\leq |B|$ is isomorphic to an object of $(\Sch/S)_{fppf}$.
Hence it suffices to observe that in the "only if" part of the proof of
Algebra, Lemma \href{https://stacks.math.columbia.edu/tag/00QO}{lemma characterize finite presentation (Tag 00QO)}
only $A$-algebras of cardinality $\leq |B|$ are used.
\end{proof}


\begin{lemma}
\label{lemma-sheafify-finite-presentation}
Let $S$ be a scheme contained in $\Sch_{fppf}$.
Let $F : (\Sch/S)_{fppf}^{opp} \to \textit{Sets}$ be a functor.
If $F$ is limit preserving then its sheafification $F^\#$ is limit preserving.
\end{lemma}

\begin{proof}
Assume $F$ is limit preserving.
It suffices to show that $F^+$ is limit preserving, since
$F^\# = (F^+)^+$, see
Sites, Theorem \href{https://stacks.math.columbia.edu/tag/00WB}{theorem plus (Tag 00WB)}.
Let $T$ be an affine scheme over $S$, and let $T = \lim T_i$ be written
as the directed limit of an inverse system of affine $S$ schemes.
Recall that $F^+(T)$ is the colimit of $\check H^0(\mathcal{V}, F)$
where the limit is over all coverings of $T$ in $(\Sch/S)_{fppf}$.
Any fppf covering of an affine scheme can be refined by a standard
fppf covering, see
Topologies, Lemma \href{https://stacks.math.columbia.edu/tag/021P}{lemma fppf affine (Tag 021P)}.
Hence we can write
$$
F^+(T)
=
\colim_{\mathcal{V}\text{ standard covering }T}
\check H^0(\mathcal{V}, F).
$$
Any $\mathcal{V} = \{T_k \to T\}_{k = 1, \ldots, n}$
in the colimit may be written as
$\mathcal{V}_i \times_{T_i} T$ for some $i$ and some standard fppf covering
$\mathcal{V}_i = \{T_{i, k} \to T_i\}_{k = 1, \ldots, n}$ of $T_i$.
Denote $\mathcal{V}_{i'} = \{T_{i', k} \to T_{i'}\}_{k = 1, \ldots, n}$
the base change for $i' \geq i$. Then we see that
\begin{align*}
\colim_{i' \geq i} \check H^0(\mathcal{V}_i, F)
& =
\colim_{i' \geq i}
\text{Equalizer}
\left(
\xymatrix{
\prod F(T_{i', k})
\ar@<1ex>[r] \ar@<-1ex>[r] &
\prod F(T_{i', k} \times_{T_{i'}} T_{i', l})
}
\right)
\\
& =
\text{Equalizer}
\left(
\xymatrix{
\colim_{i' \geq i}
\prod F(T_{i', k})
\ar@<1ex>[r] \ar@<-1ex>[r] &
\colim_{k' \geq k}
\prod F(T_{i', k} \times_{T_{i'}} T_{i', l})
}
\right)
\\
& =
\text{Equalizer}
\left(
\xymatrix{
\prod F(T_k)
\ar@<1ex>[r] \ar@<-1ex>[r] &
\prod F(T_k \times_T T_l)
}
\right)
\\
& =
\check H^0(\mathcal{V}, F)
\end{align*}
Here the second equality holds because filtered colimits are exact.
The third equality holds because $F$ is limit preserving and because
$\lim_{i' \geq i} T_{i', k} = T_k$ and
$\lim_{i' \geq i} T_{i', k} \times_{T_{i'}} T_{i', l} = T_k \times_T T_l$
by Limits, Lemma \ref{limits-lemma-scheme-over-limit}.
If we use this for all coverings at the same time we obtain
\begin{align*}
F^+(T)
& =
\colim_{\mathcal{V}\text{ standard covering }T} \check H^0(\mathcal{V}, F) \\
& =
\colim_{i \in I}
\colim_{\mathcal{V}_i\text{ standard covering }T_i}
\check H^0(T \times_{T_i}\mathcal{V}_i, F) \\
& =
\colim_{i \in I} F^+(T_i)
\end{align*}
The switch of the order of the colimits is allowed by
Categories, Lemma \href{https://stacks.math.columbia.edu/tag/002M}{lemma colimits commute (Tag 002M)}.
\end{proof}


\begin{lemma}
\label{decent-spaces-lemma-filter-quasi-compact-quasi-separated}
\begin{reference}
This result is almost identical to \href{https://stacks.math.columbia.edu/bibliography}{Bibliography: GruRay (Proposition 5.7.8)}.
\end{reference}
Let $X$ be a quasi-compact and quasi-separated algebraic space over
$\Spec(\mathbf{Z})$. There exist an integer $n$ and open subspaces
$$
\emptyset = U_{n + 1} \subset
U_n \subset U_{n - 1} \subset \ldots \subset U_1 = X
$$
with the following property: setting $T_p = U_p \setminus U_{p + 1}$
(with reduced induced subspace structure) there exists a quasi-compact
separated scheme $V_p$ and a surjective \'etale morphism $f_p : V_p \to U_p$
such that $f_p^{-1}(T_p) \to T_p$ is an isomorphism.
\end{lemma}

\begin{proof}
The proof of this lemma is identical to the proof of
Lemma \ref{decent-spaces-lemma-filter-quasi-compact}.
Observe that a quasi-separated space is reasonable, see
Lemma \ref{decent-spaces-lemma-bounded-fibres} and
Definition \ref{decent-spaces-definition-very-reasonable}.
Hence we find that $U_{n + 1} = \emptyset$ as in
Lemma \ref{decent-spaces-lemma-filter-reasonable}.
At the end of the argument we add that since $X$ is quasi-separated
the schemes $U \times_X \ldots \times_X U$ are all quasi-compact.
Hence the schemes $W_p$ are quasi-compact. Hence the
quotients $V_p = W_p/S_p$ by the symmetric group $S_p$ are quasi-compact
schemes.
\end{proof}


\begin{lemma}
\label{decent-spaces-lemma-filter-quasi-compact}
Let $S$ be a scheme. Let $X$ be a quasi-compact algebraic space
over $S$. There exist open subspaces
$$
\ldots \subset U_4 \subset U_3 \subset U_2 \subset U_1 = X
$$
with the following properties:
\begin{enumerate}
\item setting $T_p = U_p \setminus U_{p + 1}$ (with reduced induced subspace
structure) there exists a separated scheme $V_p$ and a surjective \'etale
morphism $f_p : V_p \to U_p$ such that $f_p^{-1}(T_p) \to T_p$ is an
isomorphism,
\item if $x \in |X|$ can be represented by a quasi-compact morphism
$\Spec(k) \to X$ from a field, then $x \in T_p$ for some $p$.
\end{enumerate}
\end{lemma}

\begin{proof}
By Properties of Spaces, Lemma
\href{https://stacks.math.columbia.edu/tag/03H6}{lemma quasi compact affine cover (Tag 03H6)}
we can choose an affine scheme $U$ and a surjective \'etale morphism
$U \to X$. For $p \geq 0$ set
$$
W_p = U \times_X \ldots \times_X U \setminus \text{all diagonals}
$$
where the fibre product has $p$ factors. Since $U$ is separated,
the morphism $U \to X$ is separated and all fibre products
$U \times_X \ldots \times_X U$ are separated schemes. Since $U \to X$ is
separated the diagonal $U \to U \times_X U$ is a closed immersion. Since
$U \to X$ is \'etale the diagonal $U \to U \times_X U$ is an open
immersion, see Morphisms of Spaces, Lemmas
\href{https://stacks.math.columbia.edu/tag/06CR}{lemma etale unramified (Tag 06CR)} and
\href{https://stacks.math.columbia.edu/tag/05W1}{lemma diagonal unramified morphism (Tag 05W1)}.
Similarly, all the diagonal morphisms are open and closed immersions and
$W_p$ is an open and closed subscheme of $U \times_X \ldots \times_X U$.
Moreover, the morphism
$$
U \times_X \ldots \times_X U \longrightarrow
U \times_{\Spec(\mathbf{Z})} \ldots \times_{\Spec(\mathbf{Z})} U
$$
is locally quasi-finite and separated (Morphisms of Spaces,
Lemma \href{https://stacks.math.columbia.edu/tag/03KN}{lemma fibre product after map (Tag 03KN)})
and its target is an affine scheme. Hence every finite set of points of
$U \times_X \ldots \times_X U$ is contained in an affine open, see
More on Morphisms, Lemma
\href{https://stacks.math.columbia.edu/tag/07S0}{lemma separated locally quasi finite over affine (Tag 07S0)}.
Therefore, the same is true for $W_p$.
There is a free action of the symmetric group $S_p$ on $W_p$ over $X$
(because we threw out the fix point locus from
$U \times_X \ldots \times_X U$). By the above and
Properties of Spaces, Proposition
\ref{spaces-properties-proposition-finite-flat-equivalence-global}
the quotient $V_p = W_p/S_p$ is a scheme. Since the action of
$S_p$ on $W_p$ was over $X$, there is a morphism $V_p \to X$.
Since $W_p \to X$ is \'etale and since $W_p \to V_p$ is surjective
\'etale, it follows that also $V_p \to X$ is \'etale, see
Properties of Spaces, Lemma \href{https://stacks.math.columbia.edu/tag/03FS}{lemma etale local (Tag 03FS)}.
Observe that $V_p$ is a separated scheme by
Properties of Spaces, Lemma
\href{https://stacks.math.columbia.edu/tag/0BBM}{lemma quotient separated (Tag 0BBM)}.

\medskip\noindent
We let $U_p \subset X$ be the open subspace which is the
image of $V_p \to X$. By construction a morphism $\Spec(k) \to X$ with
$k$ algebraically closed, factors through $U_p$ if and only if
$U \times_X \Spec(k)$ has $\geq p$ points; as usual observe that
$U \times_X \Spec(k)$ is scheme theoretically a disjoint union of
(possibly infinitely many) copies of $\Spec(k)$, see
Remark \href{https://stacks.math.columbia.edu/tag/03II}{remark recall (Tag 03II)}. It follows that
the $U_p$ give a filtration of $X$ as stated in the lemma.
Moreover, our morphism $\Spec(k) \to X$ factors through $T_p$
if and only if $U \times_X \Spec(k)$ has exactly $p$ points.
In this case we see that $V_p \times_X \Spec(k)$ has exactly one point.
Set $Z_p = f_p^{-1}(T_p) \subset V_p$. This is a closed subscheme of $V_p$.
Then $Z_p \to T_p$ is an \'etale morphism between
algebraic spaces which induces a bijection on $k$-valued
points for any algebraically closed field $k$. To be sure this
implies that $Z_p \to T_p$ is universally injective, whence an
open immersion by
Morphisms of Spaces, Lemma
\href{https://stacks.math.columbia.edu/tag/05W5}{lemma etale universally injective open (Tag 05W5)}
hence an isomorphism and (1) has been proved.

\medskip\noindent
Let $x : \Spec(k) \to X$ be a quasi-compact morphism where $k$ is a field.
Then the composition $\Spec(\overline{k}) \to \Spec(k) \to X$ is quasi-compact
as well (Morphisms of Spaces, Lemma
\href{https://stacks.math.columbia.edu/tag/03HG}{lemma composition quasi compact (Tag 03HG)}).
In this case the scheme $U \times_X \Spec(\overline{k})$ is
quasi-compact. In view of the fact (seen above) that it is a disjoint union
of copies of $\Spec(\overline{k})$ we find that it has finitely many points.
If the number of points is $p$, then we see that indeed $x \in T_p$ and
the proof is finished.
\end{proof}


\begin{lemma}
\label{decent-spaces-lemma-filter-reasonable}
Let $S$ be a scheme. Let $X$ be a quasi-compact, reasonable algebraic space
over $S$. There exist an integer $n$ and open subspaces
$$
\emptyset = U_{n + 1} \subset
U_n \subset U_{n - 1} \subset \ldots \subset U_1 = X
$$
with the following property: setting $T_p = U_p \setminus U_{p + 1}$
(with reduced induced subspace structure) there exists a separated scheme
$V_p$ and a surjective \'etale morphism $f_p : V_p \to U_p$ such that
$f_p^{-1}(T_p) \to T_p$ is an isomorphism.
\end{lemma}

\begin{proof}
The proof of this lemma is identical to the proof of
Lemma \ref{decent-spaces-lemma-filter-quasi-compact}.
Let $n$ be an integer bounding the degrees of
the fibres of $U \to X$ which exists as $X$ is reasonable, see
Definition \ref{decent-spaces-definition-very-reasonable}.
Then we see that $U_{n + 1} = \emptyset$ and the proof is complete.
\end{proof}


\begin{lemma}
\label{decent-spaces-lemma-bounded-fibres}
Let $S$ be a scheme. Let $X$ be an algebraic space over $S$.
Consider the following conditions on $X$:
\begin{itemize}
\item[] $(\alpha)$ For every $x \in |X|$, the equivalent conditions of
Lemma \href{https://stacks.math.columbia.edu/tag/03JS}{lemma U finite above x (Tag 03JS)}
hold.
\item[] $(\beta)$ For every $x \in |X|$, the equivalent conditions of
Lemma \href{https://stacks.math.columbia.edu/tag/03JU}{lemma R finite above x (Tag 03JU)}
hold.
\item[] $(\gamma)$ For every $x \in |X|$, the equivalent conditions of
Lemma \href{https://stacks.math.columbia.edu/tag/03JV}{lemma UR finite above x (Tag 03JV)}
hold.
\item[] $(\delta)$ The equivalent conditions of
Lemma \href{https://stacks.math.columbia.edu/tag/03JT}{lemma U universally bounded (Tag 03JT)}
hold.
\item[] $(\epsilon)$ The equivalent conditions of
Lemma \href{https://stacks.math.columbia.edu/tag/03IH}{lemma characterize very reasonable (Tag 03IH)}
hold.
\item[] $(\zeta)$ The space $X$ is Zariski locally quasi-separated.
\item[] $(\eta)$ The space $X$ is quasi-separated
\item[] $(\theta)$ The space $X$ is representable, i.e., $X$ is a scheme.
\item[] $(\iota)$ The space $X$ is a quasi-separated scheme.
\end{itemize}
We have
$$
\xymatrix{
& (\theta) \ar@{=>}[rd] & & & &  \\
(\iota) \ar@{=>}[ru] \ar@{=>}[rd] & &
(\zeta) \ar@{=>}[r] &
(\epsilon) \ar@{=>}[r] &
(\delta) \ar@{=>}[r] &
(\gamma) \ar@{<=>}[r] & (\alpha) + (\beta) \\
& (\eta) \ar@{=>}[ru] & & & &
}
$$
\end{lemma}

\begin{proof}
The implication $(\gamma) \Leftrightarrow (\alpha) + (\beta)$ is immediate.
The implications in the diamond on the left are clear from the
definitions.

\medskip\noindent
Assume $(\zeta)$, i.e., that $X$ is Zariski locally quasi-separated.
Then $(\epsilon)$ holds by
Properties of Spaces, Lemma
\href{https://stacks.math.columbia.edu/tag/03W7}{lemma quasi separated quasi compact pieces (Tag 03W7)}.

\medskip\noindent
Assume $(\epsilon)$. By
Lemma \href{https://stacks.math.columbia.edu/tag/03IH}{lemma characterize very reasonable (Tag 03IH)}
there exists
a Zariski open covering $X = \bigcup X_i$ such that for each $i$
there exists a scheme $U_i$ and a quasi-compact surjective \'etale morphism
$U_i \to X_i$. Choose an $i$ and an affine open subscheme $W \subset U_i$.
It suffices to show that $W \to X$ has universally bounded fibres, since then
the family of all these morphisms $W \to X$ covers $X$.
To do this we consider the diagram
$$
\xymatrix{
W \times_X U_i \ar[r]_-p \ar[d]_q & U_i \ar[d] \\
W \ar[r] & X
}
$$
Since $W \to X$ factors through $X_i$ we see that
$W \times_X U_i = W \times_{X_i} U_i$, and hence $q$ is quasi-compact.
Since $W$ is affine this implies that the scheme $W \times_X U_i$
is quasi-compact. Thus we may apply
Morphisms, Lemma
\href{https://stacks.math.columbia.edu/tag/03JA}{lemma locally quasi finite qc source universally bounded (Tag 03JA)}
and we conclude that $p$ has universally bounded fibres. From
Lemma \href{https://stacks.math.columbia.edu/tag/03JO}{lemma descent universally bounded (Tag 03JO)}
we conclude that $W \to X$ has universally bounded fibres as well.

\medskip\noindent
Assume $(\delta)$. Let $U$ be an affine scheme, and let $U \to X$ be an \'etale
morphism. By assumption the fibres of the morphism $U \to X$ are universally
bounded. Thus also the fibres of both projections $R = U \times_X U \to U$
are universally bounded, see
Lemma \href{https://stacks.math.columbia.edu/tag/03JN}{lemma base change universally bounded (Tag 03JN)}.
And by
Lemma \href{https://stacks.math.columbia.edu/tag/03JM}{lemma composition universally bounded (Tag 03JM)}
also the fibres of $R \to X$ are universally bounded.
Hence for any $x \in X$ the fibres of $|U| \to |X|$ and $|R| \to |X|$
over $x$ are finite, see
Lemma \href{https://stacks.math.columbia.edu/tag/03JQ}{lemma universally bounded finite fibres (Tag 03JQ)}.
In other words, the equivalent conditions of
Lemma \href{https://stacks.math.columbia.edu/tag/03JV}{lemma UR finite above x (Tag 03JV)}
hold. This proves that $(\delta) \Rightarrow (\gamma)$.
\end{proof}


\begin{proposition}
\label{spaces-properties-proposition-finite-flat-equivalence-global}
Let $S$ be a scheme.
Let $(U, R, s, t, c)$ be a groupoid scheme over $S$.
Assume
\begin{enumerate}
\item $s, t : R \to U$ finite locally free,
\item $j = (t, s)$ is an equivalence relation, and
\item every nonempty closed subset $Z$ of $U$ contains a point
$u$ whose $R$-equivalence class $t(s^{-1}(\{u\}))$ is contained
in an affine open of $U$\footnote{Let $E \subset U$ be the
set of points $u$ such that $t(s^{-1}(\{u\}))$
is contained in an affine open of $U$. Condition (3) holds
if $E = U$, or if every finite type point of $U$ is in $E$, or
if every $u \in U$ specializes to a point of $E$.}.
\end{enumerate}
Then there exists a finite locally free morphism $U \to M$
of schemes over $S$ such that $R = U \times_M U$ and such that $M$
represents the quotient sheaf $U/R$ in the fppf topology.
\end{proposition}

\begin{proof}
By assumption (3) and
Groupoids, Lemma \ref{groupoids-lemma-find-invariant-affine}
we can find an open covering $U = \bigcup U_i$ such that each $U_i$
is an $R$-invariant affine open of $U$. Set $R_i = R|_{U_i}$.
Consider the fppf sheaves $F = U/R$ and $F_i = U_i/R_i$.
By Spaces, Lemma \href{https://stacks.math.columbia.edu/tag/02WU}{lemma finding opens (Tag 02WU)} the morphisms
$F_i \to F$ are representable and open immersions.
By Groupoids, Proposition \ref{groupoids-proposition-finite-flat-equivalence}
the sheaves $F_i$ are representable by affine schemes.
If $T$ is a scheme and $T \to F$ is a morphism, then $V_i = F_i \times_F T$
is open in $T$ and we claim that $T = \bigcup V_i$. Namely,
fppf locally on $T$ we can lift $T \to F$ to a morphism
$f : T \to U$ and in that case $f^{-1}(U_i) \subset V_i$.
Hence we conclude that $F$ is representable by a scheme, see
Schemes, Lemma \href{https://stacks.math.columbia.edu/tag/01JJ}{lemma glue functors (Tag 01JJ)}.
\end{proof}


\begin{lemma}
\label{groupoids-lemma-find-invariant-affine}
Let $S$ be a scheme.
Let $(U, R, s, t, c)$ be a groupoid scheme over $S$.
Assume $s$, $t$ are finite locally free.
Let $u \in U$ be a point such that $t(s^{-1}(\{u\}))$
is contained in an affine open of $U$.
Then there exists an $R$-invariant affine open neighbourhood
of $u$ in $U$.
\end{lemma}

\begin{proof}
Since $s$ is finite locally free it has finite fibres. Hence
$t(s^{-1}(\{u\})) = \{u_1, \ldots, u_n\}$ is a finite set.
Note that $u \in \{u_1, \ldots, u_n\}$.
Let $W \subset U$ be an affine open containing $\{u_1, \ldots, u_n\}$,
in particular $u \in W$. Consider
$Z = R \setminus s^{-1}(W) \cap t^{-1}(W)$. This is a closed subset
of $R$. The image $t(Z)$ is a closed subset of $U$ which can be loosely
described as the set of points of $U$ which are $R$-equivalent to a point
of $U \setminus W$. Hence $W' = U \setminus t(Z)$ is an $R$-invariant, open
subscheme of $U$ contained in $W$, and $\{u_1, \ldots, u_n\} \subset W'$.
Picture
$$
\{u_1, \ldots, u_n\} \subset W' \subset W \subset U.
$$
Let $f \in \Gamma(W, \mathcal{O}_W)$ be an element such that
$\{u_1, \ldots, u_n\} \subset D(f) \subset W'$. Such an $f$ exists by
Algebra, Lemma \href{https://stacks.math.columbia.edu/tag/00DS}{lemma silly (Tag 00DS)}. By our choice of $W'$ we
have $s^{-1}(W') \subset t^{-1}(W)$, and hence we get a diagram
$$
\xymatrix{
s^{-1}(W') \ar[d]_s \ar[r]_-t & W \\
W'
}
$$
The vertical arrow is finite locally free by assumption. Set
$$
g = \text{Norm}_s(t^\sharp f) \in \Gamma(W', \mathcal{O}_{W'})
$$
By construction $g$ is a function on $W'$ which is
nonzero in $u$, as $t^\sharp(f)$ is nonzero in each of the points of
$R$ lying over $u$, since $f$ is nonzero in $u_1, \ldots, u_n$.
Similarly, $D(g) \subset W'$ is equal to the
set of points $w$ such that $f$ is not zero in any of the points
equivalent to $w$. This means that $D(g)$ is an
$R$-invariant affine open of $W'$. The final picture is
$$
\{u_1, \ldots, u_n\} \subset D(g) \subset D(f) \subset W' \subset W \subset U
$$
and hence we win.
\end{proof}


\begin{proposition}
\label{groupoids-proposition-finite-flat-equivalence}
Let $S$ be a scheme.
Let $(U, R, s, t, c)$ be a groupoid scheme over $S$.
Assume
\begin{enumerate}
\item $U = \Spec(A)$, and $R = \Spec(B)$ are affine,
\item $s, t : R \to U$ finite locally free, and
\item $j = (t, s)$ is an equivalence relation.
\end{enumerate}
In this case, let $C \subset A$ be as in
(\ref{groupoids-equation-invariants}). Then $U \to M = \Spec(C)$
is finite locally free and $R = U \times_M U$.
Moreover, $M$ represents the quotient sheaf $U/R$
in the fppf topology (see Definition \ref{groupoids-definition-quotient-sheaf}).
\end{proposition}

\begin{proof}
During this proof we use the notation $s, t : A \to B$
instead of the notation $s^\sharp, t^\sharp$.
By Lemma \ref{groupoids-lemma-criterion-quotient-representable}
it suffices to show that $C \to A$ is finite locally free
and that the map
$$
t \otimes s : A \otimes_C A \longrightarrow B
$$
is an isomorphism. First, note that $j$ is a monomorphism, and
also finite (since already $s$ and $t$ are finite). Hence we see
that $j$ is a closed immersion by
Morphisms, Lemma \href{https://stacks.math.columbia.edu/tag/03BB}{lemma finite monomorphism closed (Tag 03BB)}.
Hence $A \otimes_C A \to B$ is surjective.

\medskip\noindent
We will perform base change by flat ring maps $C \to C'$ as in
Lemma \ref{groupoids-lemma-invariants-base-change}, and we will use that
formation of invariants commutes with flat base change, see
part (3) of the lemma cited.
We will show below that for every prime $\mathfrak p \subset C$, there exists
a local flat ring map $C_{\mathfrak p} \to C_{\mathfrak p}'$
such that the result holds after a base change to $C_{\mathfrak p}'$.
This implies immediately
that $A \otimes_C A \to B$ is injective (use
Algebra, Lemma \href{https://stacks.math.columbia.edu/tag/00HN}{lemma characterize zero local (Tag 00HN)}).
It also implies that $C \to A$ is flat, by combining
Algebra, Lemmas \href{https://stacks.math.columbia.edu/tag/00HR}{lemma local flat ff (Tag 00HR)},
\href{https://stacks.math.columbia.edu/tag/00HT}{lemma flat localization (Tag 00HT)}, and
\href{https://stacks.math.columbia.edu/tag/00HJ}{lemma flatness descends (Tag 00HJ)}. Then since $U \to \Spec(C)$
is surjective also (Lemma \ref{groupoids-lemma-points}) we conclude that $C \to A$
is faithfully flat. Then the isomorphism $B \cong A \otimes_C A$
implies that $A$ is a finitely presented $C$-module, see
Algebra, Lemma \href{https://stacks.math.columbia.edu/tag/03C4}{lemma descend properties modules (Tag 03C4)}.
Hence $A$ is finite locally free over $C$, see
Algebra, Lemma \href{https://stacks.math.columbia.edu/tag/00NX}{lemma finite projective (Tag 00NX)}.

\medskip\noindent
By Lemma \ref{groupoids-lemma-finite-locally-free-disjoint-free}
we know that $A$ is a finite
product of rings $A_r$ and $B$ is a finite product of rings $B_r$
such that the groupoid scheme decomposes accordingly (see the proof
of Lemma \ref{groupoids-lemma-integral-over-invariants}).
Then also $C$ is a product of rings $C_r$ and
correspondingly $C'$ decomposes as a product. Hence we may and do
assume that the ring maps $s, t : A \to B$ are finite
locally free of a fixed rank $r$.

\medskip\noindent
The local ring maps $C_{\mathfrak p} \to C_{\mathfrak p}'$ we are going
to use are any local flat ring maps such that the residue field of
$C_{\mathfrak p}'$ is infinite.
By Algebra, Lemma \ref{algebra-lemma-flat-local-given-residue-field}
such local ring maps exist.

\medskip\noindent
Assume $C$ is a local ring with maximal ideal $\mathfrak m$ and
infinite residue field, and assume that $s, t : A \to B$ is
finite locally free of constant rank $r > 0$.
Since $C \subset A$ is integral (Lemma \ref{groupoids-lemma-integral-over-invariants})
all primes lying over $\mathfrak m$ are maximal, and all maximal
ideals of $A$ lie over $\mathfrak m$. Similarly for $C \subset B$.
Pick a maximal ideal $\mathfrak m'$
of $A$ lying over $\mathfrak m$ (exists by Lemma \ref{groupoids-lemma-points}).
Since $t : A \to B$ is finite locally free there exist at most finitely
many maximal ideals of $B$ lying over $\mathfrak m'$. Hence we conclude
(by Lemma \ref{groupoids-lemma-points} again)
that $A$ has finitely many maximal ideals, i.e.,
$A$ is semi-local. This in turn implies that $B$ is semi-local as
well. OK, and now, because $t \otimes s : A \otimes_C A \to B$ is surjective,
we can apply
Algebra, Lemma \ref{algebra-lemma-semi-local-module-basis-in-submodule}
to the ring map $C \to A$, the $A$-module $M = B$ (seen as an $A$-module
via $t$) and the $C$-submodule $s(A) \subset B$. This lemma implies that there
exist $x_1, \ldots, x_r \in A$ such that $M$ is free over $A$
on the basis $s(x_1), \ldots, s(x_r)$. Hence we conclude that $C \to A$
is finite free and $B \cong A \otimes_C A$ by applying
Lemma \ref{groupoids-lemma-basis}.
\end{proof}


\begin{theorem}
\label{spaces-theorem-presentation}
Let $S$ be a scheme. Let $U$ be a scheme over $S$.
Let $j = (s, t) : R \to U \times_S U$
be an \'etale equivalence relation on $U$ over $S$.
Then the quotient $U/R$ is an algebraic space,
and $U \to U/R$ is \'etale and surjective, in other words
$(U, R, U \to U/R)$ is a presentation of $U/R$.
\end{theorem}

\begin{proof}
By Lemma \href{https://stacks.math.columbia.edu/tag/02WV}{lemma when it works it works (Tag 02WV)}
it suffices to prove that $U/R$ is an algebraic space.
Let $U' \to U$ be a surjective, \'etale morphism.
Then $\{U' \to U\}$ is in particular an fppf covering.
Let $R'$ be the restriction of $R$ to $U'$, see
Groupoids, Definition \ref{groupoids-definition-restrict-relation}.
According to
Groupoids, Lemma \href{https://stacks.math.columbia.edu/tag/02VH}{lemma quotient groupoid restrict (Tag 02VH)}
we see that $U/R \cong U'/R'$.
By Lemma \href{https://stacks.math.columbia.edu/tag/02WT}{lemma pullback etale equivalence relation (Tag 02WT)} $R'$ is an
\'etale equivalence relation on $U'$. Thus we may replace $U$ by $U'$.

\medskip\noindent
We apply the previous remark to $U' = \coprod U_i$, where
$U = \bigcup U_i$ is an affine open covering of $U$. Hence we
may and do assume that $U = \coprod U_i$ where
each $U_i$ is an affine scheme.

\medskip\noindent
Consider the restriction $R_i$ of $R$ to $U_i$.
By Lemma \href{https://stacks.math.columbia.edu/tag/02WT}{lemma pullback etale equivalence relation (Tag 02WT)}
this is an \'etale equivalence relation.
Set $F_i = U_i/R_i$ and $F = U/R$.
It is clear that $\coprod F_i \to F$ is surjective.
By Lemma \href{https://stacks.math.columbia.edu/tag/02WU}{lemma finding opens (Tag 02WU)} each $F_i \to F$
is representable, and an open immersion.
By Lemma \ref{spaces-lemma-presentation-quasi-compact}
applied to $(U_i, R_i)$ we see that $F_i$ is an algebraic space.
Then by Lemma \href{https://stacks.math.columbia.edu/tag/02WV}{lemma when it works it works (Tag 02WV)} we see that
$U_i \to F_i$ is \'etale and surjective.
From Lemma \href{https://stacks.math.columbia.edu/tag/02WQ}{lemma coproduct algebraic spaces (Tag 02WQ)}
it follows that $\coprod F_i$ is an algebraic space.
Finally, we have verified all
hypotheses of Lemma \href{https://stacks.math.columbia.edu/tag/02WR}{lemma glueing algebraic spaces (Tag 02WR)}
and it follows that $F = U/R$ is an algebraic space.
\end{proof}


\begin{lemma}
\label{spaces-lemma-presentation-quasi-compact}
Let $S$ be a scheme.
Let $U$ be a scheme over $S$.
Let $j = (s, t) : R \to U \times_S U$
be an \'etale equivalence relation on $U$ over $S$.
Assume that $U$ is affine. Then the quotient $F = U/R$
is an algebraic space, and $U \to F$ is \'etale and surjective.
\end{lemma}

\begin{proof}
Since $j : R \to U \times_S U$ is a monomorphism we see that $j$ is separated
(see Schemes, Lemma \href{https://stacks.math.columbia.edu/tag/01L4}{lemma monomorphism separated (Tag 01L4)}).
Since $U$ is affine we see that $U \times_S U$
(which comes equipped with a monomorphism into the affine scheme
$U \times U$) is separated. Hence we see that $R$ is separated.
In particular the morphisms $s, t$ are separated as well as \'etale.

\medskip\noindent
Since the composition $R \to U \times_S U \to U$ is
locally of finite type we conclude that
$j$ is locally of finite type (see
Morphisms, Lemma \href{https://stacks.math.columbia.edu/tag/01T8}{lemma permanence finite type (Tag 01T8)}).
As $j$ is also a monomorphism it has finite fibres and
we see that $j$ is locally quasi-finite by
Morphisms, Lemma \href{https://stacks.math.columbia.edu/tag/02NG}{lemma finite fibre (Tag 02NG)}.
Altogether we see that $j$ is separated and locally quasi-finite.

\medskip\noindent
Our first step is to show that the quotient map
$c : U \to F$ is representable.
Consider a scheme $T$ and a morphism $a : T \to F$.
We have to show that the sheaf $G = T \times_{a, F, c} U$
is representable.
As seen in the proofs of Lemmas \href{https://stacks.math.columbia.edu/tag/02WU}{lemma finding opens (Tag 02WU)} and
\href{https://stacks.math.columbia.edu/tag/02WV}{lemma when it works it works (Tag 02WV)} there exists an fppf covering
$\{\varphi_i : T_i \to T\}_{i \in I}$ and morphisms $a_i : T_i \to U$
such that $a_i \times a_{i'} : T_i \times_T T_{i'} \to U \times_S U$
factors through $R$, and such that $c \circ a_i = a \circ \varphi_i$.
As in the proof of Lemma \href{https://stacks.math.columbia.edu/tag/02WV}{lemma when it works it works (Tag 02WV)} we see that
\begin{eqnarray*}
T_i \times_{\varphi_i, T} G & = &
T_i \times_{\varphi_i, T} T \times_{a, U/R, c} U \\
& = & T_i \times_{c \circ a_i, U/R, c} U \\
& = & T_i \times_{a_i, U} U \times_{c, U/R, c} U \\
& = & T_i \times_{a_i, U, t} R
\end{eqnarray*}
Since $t$ is separated and \'etale, and in particular
separated and locally quasi-finite (by Morphisms, Lemmas
\href{https://stacks.math.columbia.edu/tag/02V5}{lemma unramified quasi finite (Tag 02V5)} and
\href{https://stacks.math.columbia.edu/tag/02GV}{lemma flat unramified etale (Tag 02GV)})
we see that the restriction
of $G$ to each $T_i$ is representable by a morphism of schemes
$X_i \to T_i$ which is separated and locally quasi-finite. By
Descent, Lemma \href{https://stacks.math.columbia.edu/tag/02W5}{lemma descent data sheaves (Tag 02W5)}
we obtain a descent datum $(X_i, \varphi_{ii'})$ relative
to the fppf-covering $\{T_i \to T\}$. Since each
$X_i \to T_i$ is separated and locally quasi-finite we see by
More on Morphisms, Lemma
\href{https://stacks.math.columbia.edu/tag/02W8}{lemma separated locally quasi finite morphisms fppf descend (Tag 02W8)}
that this descent datum is effective.
Hence by
Descent, Lemma \href{https://stacks.math.columbia.edu/tag/02W5}{lemma descent data sheaves (Tag 02W5)} (2)
we conclude that $G$ is representable as desired.

\medskip\noindent
The second step of the proof is to show that $U \to F$ is surjective and
\'etale. This is clear from the above since in the first step above we
saw that $G = T \times_{a, F, c} U$ is a scheme over $T$ which base changes
to schemes $X_i \to T_i$ which are surjective and \'etale. Thus $G \to T$
is surjective and \'etale (see
Remark \href{https://stacks.math.columbia.edu/tag/02WH}{remark list properties fpqc local base (Tag 02WH)}).
Alternatively one can reread the proof of
Lemma \href{https://stacks.math.columbia.edu/tag/02WV}{lemma when it works it works (Tag 02WV)} in the current
situation.

\medskip\noindent
The third and final step is to show that the diagonal map $F \to F \times F$
is representable. We first observe that the diagram
$$
\xymatrix{
R \ar[r] \ar[d]_j & F \ar[d]^\Delta \\
U \times_S U \ar[r] & F \times F
}
$$
is a fibre product square. By
Lemma \href{https://stacks.math.columbia.edu/tag/02WC}{lemma product representable transformations (Tag 02WC)} the morphism
$U \times_S U \to F \times F$ is representable (note that
$h_U \times h_U = h_{U \times_S U}$). Moreover, by
Lemma \href{https://stacks.math.columbia.edu/tag/02WM}{lemma product representable transformations property (Tag 02WM)}
the morphism $U \times_S U \to F \times F$ is surjective
and \'etale (note also that \'etale and surjective occur in the lists of
Remarks \href{https://stacks.math.columbia.edu/tag/02WH}{remark list properties fpqc local base (Tag 02WH)}
and \href{https://stacks.math.columbia.edu/tag/02WG}{remark list properties stable composition (Tag 02WG)}).
It follows either from
Lemma \href{https://stacks.math.columbia.edu/tag/02WB}{lemma base change representable transformations (Tag 02WB)}
and the diagram above, or by writing $R \to F$ as $R \to U \to F$ and
Lemmas
\href{https://stacks.math.columbia.edu/tag/02W9}{lemma morphism schemes gives representable transformation (Tag 02W9)} and
\href{https://stacks.math.columbia.edu/tag/02WA}{lemma composition representable transformations (Tag 02WA)} that
$R \to F$ is representable as well. Let $T$ be a scheme and let
$a : T \to F \times F$ be a morphism. We have to show that
$G = T \times_{a, F \times F, \Delta} F$ is representable.
By what was said above the morphism (of schemes)
$$
T' = (U \times_S U) \times_{F \times F, a} T \longrightarrow T
$$
is surjective and \'etale. Hence $\{T' \to T\}$ is an \'etale
covering of $T$. Note also that
$$
T' \times_T G = T' \times_{U \times_S U, j} R
$$
as can be seen contemplating the following cube
$$
\xymatrix{
& R \ar[rr] \ar[dd] & & F \ar[dd] \\
T' \times_T G \ar[rr] \ar[dd] \ar[ru] & & G \ar[dd] \ar[ru] & \\
& U \times_S U \ar'[r][rr] & & F \times F \\
T' \ar[rr] \ar[ru] & & T \ar[ru]
}
$$
Hence we see that the restriction of $G$ to $T'$ is representable
by a scheme $X$, and moreover that the morphism $X \to T'$ is
a base change of the morphism $j$. Hence $X \to T'$ is
separated and locally quasi-finite (see second paragraph of the proof).
By Descent, Lemma \href{https://stacks.math.columbia.edu/tag/02W5}{lemma descent data sheaves (Tag 02W5)}
we obtain a descent datum $(X, \varphi)$ relative
to the fppf-covering $\{T' \to T\}$. Since
$X \to T'$ is separated and locally quasi-finite we see by
More on Morphisms, Lemma
\href{https://stacks.math.columbia.edu/tag/02W8}{lemma separated locally quasi finite morphisms fppf descend (Tag 02W8)}
that this descent datum is effective.
Hence by
Descent, Lemma \href{https://stacks.math.columbia.edu/tag/02W5}{lemma descent data sheaves (Tag 02W5)} (2)
we conclude that $G$ is representable as desired.
\end{proof}


\begin{lemma}
\label{groupoids-lemma-criterion-fibre-product}
Let $S$ be a scheme. Let $f : (U, R, j) \to (U', R', j')$ be a morphism
between equivalence relations over $S$. Assume that
$$
\xymatrix{
R \ar[d]_s \ar[r]_f & R' \ar[d]^{s'} \\
U \ar[r]^f & U'
}
$$
is cartesian. For any
$\tau \in  \{Zariski, \etale, fppf, smooth, syntomic\}$
the diagram
$$
\xymatrix{
U \ar[d] \ar[r] & U/R \ar[d]^f \\
U' \ar[r] & U'/R'
}
$$
is a fibre product square of $\tau$-sheaves.
\end{lemma}

\begin{proof}
By Lemma \href{https://stacks.math.columbia.edu/tag/045Y}{lemma quotient pre equivalence (Tag 045Y)} the quotient sheaves
have a simple description which we will use below without further mention.
We first show that
$$
U \longrightarrow U' \times_{U'/R'} U/R
$$
is injective. Namely, assume $a, b \in U(T)$ map to the same element
on the right hand side. Then $f(a) = f(b)$. After replacing $T$ by the
members of a $\tau$-covering we may assume that there exists an
$r \in R(T)$ such that $a = s(r)$ and $b = t(r)$. Then $r' = f(r)$
is a $T$-valued point of $R'$ with $s'(r') = t'(r')$. Hence
$r' = e'(f(a))$ (where $e'$ is the identity of the groupoid
scheme associated to $j'$, see Lemma \href{https://stacks.math.columbia.edu/tag/0233}{lemma equivalence groupoid (Tag 0233)}).
Because the first diagram of the lemma is cartesian this implies
that $r$ has to equal $e(a)$. Thus $a = b$.

\medskip\noindent
Finally, we show that the displayed arrow is surjective. Let
$T$ be a scheme over $S$ and let $(a', \overline{b})$ be a section
of the sheaf $U' \times_{U'/R'} U/R$ over $T$. After replacing $T$
by the members of a $\tau$-covering we may assume that $\overline{b}$
is the class of an element $b \in U(T)$. After replacing $T$
by the members of a $\tau$-covering we may assume that there exists
an $r' \in R'(T)$ such that $a' = t(r')$ and $s'(r') = f(b)$.
Because the first diagram of the lemma is cartesian we can find
$r \in R(T)$ such that $s(r) = b$ and $f(r) = r'$. Then it is clear
that $a = t(r) \in U(T)$ is a section which maps to
$(a', \overline{b})$.
\end{proof}


\begin{lemma}
\label{groupoids-lemma-determinant-trick}
Let $S$ be a scheme. Let $(U, R, s, t, c)$ be a groupoid scheme over $S$.
Assume $U = \Spec(A)$ and $R = \Spec(B)$ are affine and
$s, t : R \to U$ finite locally free.
Let $C$ be as in (\ref{groupoids-equation-invariants}).
Let $f \in A$. Then $\text{Norm}_{s^\sharp}(t^\sharp(f)) \in C$.
\end{lemma}

\begin{proof}
Consider the commutative diagram
$$
\xymatrix{
& U & \\
R \ar[d]_s \ar[ru]^t &
R \times_{s, U, t} R
\ar[l]^-{\text{pr}_0} \ar[d]^{\text{pr}_1} \ar[r]_-c &
R \ar[d]^s \ar[lu]_t \\
U & R \ar[l]_t \ar[r]^s & U
}
$$
of Lemma \ref{groupoids-lemma-diagram}.
Think of $f \in \Gamma(U, \mathcal{O}_U)$. The commutativity of the
top part of the diagram shows that
$\text{pr}_0^\sharp(t^\sharp(f)) = c^\sharp(t^\sharp(f))$ as elements of
$\Gamma(R \times_{S, U, t} R, \mathcal{O})$.
Looking at the right lower cartesian square
the compatibility of the norm construction with base change shows that
$s^\sharp(\text{Norm}_{s^\sharp}(t^\sharp(f))) =
\text{Norm}_{\text{pr}_1^\sharp}(c^\sharp(t^\sharp(f)))$.
Similarly we get
$t^\sharp(\text{Norm}_{s^\sharp}(t^\sharp(f))) =
\text{Norm}_{\text{pr}_1^\sharp}(\text{pr}_0^\sharp(t^\sharp(f)))$.
Hence by the first equality of this proof we see that
$s^\sharp(\text{Norm}_{s^\sharp}(t^\sharp(f))) =
t^\sharp(\text{Norm}_{s^\sharp}(t^\sharp(f)))$ as desired.
\end{proof}


\begin{lemma}
\label{groupoids-lemma-finite-locally-free-disjoint-free}
Let $S$ be a scheme. Let $(U, R, s, t, c)$ be a groupoid scheme over $S$.
Assume $s, t : R \to U$ finite locally free.
Then
$$
U = \coprod\nolimits_{r \geq 1} U_r
$$
is a disjoint union of $R$-invariant opens such that the restriction $R_r$ of
$R$ to $U_r$ has the property that $s, t : R_r \to U_r$ are finite locally
free of rank $r$.
\end{lemma}

\begin{proof}
By
Morphisms, Lemma \href{https://stacks.math.columbia.edu/tag/04MH}{lemma finite locally free (Tag 04MH)}
there exists a decomposition
$U = \coprod\nolimits_{r \geq 0} U_r$
such that $s : s^{-1}(U_r) \to U_r$ is finite locally free of rank $r$.
As $s$ is surjective we see that $U_0 = \emptyset$.
Note that $u \in U_r \Leftrightarrow$ if and only if the scheme theoretic fibre
$s^{-1}(u)$ has degree $r$ over $\kappa(u)$. Now, if $z \in R$ with $s(z) = u$
and $t(z) = u'$ then using notation as in Lemma \ref{groupoids-lemma-diagram}
$$
\text{pr}_1^{-1}(z) \to \Spec(\kappa(z))
$$
is the base change of both
$s^{-1}(u) \to \Spec(\kappa(u))$ and $s^{-1}(u') \to \Spec(\kappa(u'))$
by the lemma cited. Hence $u \in U_r \Leftrightarrow u' \in U_r$,
in other words, the open subsets $U_r$ are $R$-invariant.
In particular the restriction of $R$ to $U_r$ is just
$s^{-1}(U_r)$ and $s : R_r \to U_r$ is finite locally free of rank $r$.
As $t : R_r \to U_r$ is isomorphic to $s$ by the inverse of $R_r$
we see that it has also rank $r$.
\end{proof}


\begin{lemma}
\label{groupoids-lemma-integral-over-invariants}
Let $S$ be a scheme. Let $(U, R, s, t, c)$ be a groupoid scheme over $S$.
Assume $U = \Spec(A)$ and $R = \Spec(B)$ are affine and
$s, t : R \to U$ finite locally free.
Let $C \subset A$ be as in (\ref{groupoids-equation-invariants}).
Then $A$ is integral over $C$.
\end{lemma}

\begin{proof}
First, by Lemma \ref{groupoids-lemma-finite-locally-free-disjoint-free}
we know that $(U, R, s, t, c)$ is a disjoint union
of groupoid schemes $(U_r, R_r, s, t, c)$ such that each $s, t : R_r \to U_r$
has constant rank $r$. As $U$ is quasi-compact, we have $U_r = \emptyset$ for
almost all $r$. It suffices to prove the lemma for each $(U_r, R_r, s, t, c)$
and hence we may assume that $s, t$ are finite locally free of rank $r$.

\medskip\noindent
Assume that $s, t$ are finite locally free of rank $r$.
Let $f \in A$. Consider the element $x - f \in A[x]$, where we think
of $x$ as the coordinate on $\mathbf{A}^1$.
Since
$$
(U \times \mathbf{A}^1, R \times \mathbf{A}^1,
s \times \text{id}_{\mathbf{A}^1},
t \times \text{id}_{\mathbf{A}^1},
c \times \text{id}_{\mathbf{A}^1})
$$
is also a groupoid scheme with finite source and target, we may apply
Lemma \ref{groupoids-lemma-determinant-trick} to it and we see that
$P(x) = \text{Norm}_{s^\sharp}(t^\sharp(x - f))$
is an element of $C[x]$. Because $s^\sharp : A \to B$ is finite locally
free of rank $r$ we see that $P$ is monic of degree $r$.
Moreover $P(f) = 0$ by Cayley-Hamilton
(Algebra, Lemma \href{https://stacks.math.columbia.edu/tag/00DX}{lemma charpoly (Tag 00DX)}).
More precisely, Cayley-Hamilton says the polynomial
$s^\sharp(P) \in B[x]$ annihilates the element $t^\sharp(f)$.
Since the coefficients of $P$ are in $C$ we get that
$t^\sharp(P(f)) = 0$ in $B$. Since $t$ is faithfully flat,
we get $P(f) = 0$.
\end{proof}


\begin{lemma}
\label{groupoids-lemma-invariants-base-change}
Let $S$ be a scheme. Let $(U, R, s, t, c)$ be a groupoid scheme over $S$.
Assume $U = \Spec(A)$ and $R = \Spec(B)$ are affine and
$s, t : R \to U$ finite locally free. Let $C \subset A$ be as in
(\ref{groupoids-equation-invariants}). Let $C \to C'$ be a ring map, and set
$U' = \Spec(A \otimes_C C')$,
$R' = \Spec(B \otimes_C C')$.
Then
\begin{enumerate}
\item The maps $s, t, c$ induce maps $s', t', c'$ such that
$(U', R', s', t', c')$ is a groupoid scheme. Let $C^1 \subset A'$
be the $R'$-invariant functions on $U'$.
\item The canonical map $\varphi : C' \to C^1$ satisfies
\begin{enumerate}
\item for every $f \in C^1$ there exists an $n > 0$ and a
polynomial $P \in C'[x]$ whose image in $C^1[x]$ is
$(x - f)^n$, and
\item for every $f \in \Ker(\varphi)$ there exists
an $n > 0$ such that $f^n = 0$.
\end{enumerate}
\item If $C \to C'$ is flat then $\varphi$ is an isomorphism.
\end{enumerate}
\end{lemma}

\begin{proof}
The proof of part (1) is omitted. Let us denote $A' = A \otimes_C C'$ and
$B' = B \otimes_C C'$. Then we have
$$
C^1
= \{a \in A' \mid (t')^\sharp(a) = (s')^\sharp(a) \}
= \{a \in A \otimes_C C' \mid t^\sharp \otimes 1(a) = s^\sharp \otimes 1(a) \}.
$$
In other words, $C^1$ is the kernel of the difference map
$(t^\sharp - s^\sharp) \otimes 1$ which is just the base change
of the $C$-linear map $t^\sharp - s^\sharp : A \to B$ by $C \to C'$.
Hence (3) follows.

\medskip\noindent
Proof of part (2)(b). Since $C \to A$ is integral
(Lemma \ref{groupoids-lemma-integral-over-invariants}) and injective we see that
$\Spec(A) \to \Spec(C)$ is surjective, see
Algebra, Lemma \href{https://stacks.math.columbia.edu/tag/00GQ}{lemma integral overring surjective (Tag 00GQ)}.
Thus also $\Spec(A') \to \Spec(C')$ is surjective
as a base change of a surjective morphism
(Morphisms, Lemma \href{https://stacks.math.columbia.edu/tag/01S1}{lemma base change surjective (Tag 01S1)}).
Hence $\Spec(C^1) \to \Spec(C')$ is surjective also.
This implies (2)(b) holds for example by
Algebra, Lemma \href{https://stacks.math.columbia.edu/tag/00FL}{lemma image dense generic points (Tag 00FL)}.

\medskip\noindent
Proof of part (2)(a). By Lemma \ref{groupoids-lemma-finite-locally-free-disjoint-free}
our groupoid scheme $(U, R, s, t, c)$ decomposes as a finite disjoint union
of groupoid schemes $(U_r, R_r, s, t, c)$ such that $s, t : R_r \to U_r$
are finite locally free of rank $r$. Pulling back by $U' = \Spec(C') \to U$
we obtain a similar decomposition of $U'$ and $U^1 = \Spec(C^1)$.
We will show in the next paragraph that (2)(a) holds for the corresponding
system of rings $A_r, B_r, C_r, C'_r, C^1_r$ with $n = r$.
Then given $f \in C^1$ let $P_r \in C_r[x]$ be the polynomial
whose image in $C^1_r[x]$ is the image of $(x - f)^r$.
Choosing a sufficiently divisible integer $n$ we see that
there is a polynomial $P \in C'[x]$ whose image in $C^1[x]$ is
$(x - f)^n$; namely, we take $P$ to be the unique element of
$C'[x]$ whose image in $C'_r[x]$ is $P_r^{n/r}$.

\medskip\noindent
In this paragraph we prove (2)(a) in case the ring maps
$s^\sharp, t^\sharp : A \to B$ are finite locally free of a fixed rank $r$.
Let $f \in C^1 \subset A' = A \otimes_C C'$. Choose a flat
$C$-algebra $D$ and a surjection $D \to C'$. Choose a lift
$g \in A \otimes_C D$ of $f$.
Consider the polynomial
$$
P = \text{Norm}_{s^\sharp \otimes 1}(t^\sharp \otimes 1(x - g))
$$
in $(A \otimes_C D)[x]$. By Lemma \ref{groupoids-lemma-determinant-trick}
and part (3) of the current lemma the coefficients of $P$ are in $D$
(compare with the proof of Lemma \ref{groupoids-lemma-integral-over-invariants}).
On the other hand, the image of $P$ in $(A \otimes_C C')[x]$ is
$(x - f)^r$ because $t^\sharp \otimes 1(x - f) = s^\sharp(x - f)$
and $s^\sharp$ is finite locally free of rank $r$.
This proves what we want with $P$ as in the statement (2)(a)
given by the image of our $P$ under the map $D[x] \to C'[x]$.
\end{proof}


\begin{lemma}
\label{groupoids-lemma-points}
Let $S$ be a scheme. Let $(U, R, s, t, c)$ be a groupoid scheme over $S$.
Assume $U = \Spec(A)$ and $R = \Spec(B)$ are affine and
$s, t : R \to U$ finite locally free. Let $C \subset A$ be as in
(\ref{groupoids-equation-invariants}). Then $U \to M = \Spec(C)$ has
the following properties:
\begin{enumerate}
\item the map on points $|U| \to |M|$ is surjective and
$u_0, u_1 \in |U|$ map to the same point if and only if
there exists a $r \in |R|$ with $t(r) = u_0$ and $s(r) = u_1$, in
a formula
$$
|M| = |U|/|R|
$$
\item for any algebraically closed field $k$ we have
$$
M(k) = U(k)/R(k)
$$
\end{enumerate}
\end{lemma}

\begin{proof}
Since $C \to A$ is integral (Lemma \ref{groupoids-lemma-integral-over-invariants})
and injective we see that $\Spec(A) \to \Spec(C)$ is surjective, see
Algebra, Lemma \href{https://stacks.math.columbia.edu/tag/00GQ}{lemma integral overring surjective (Tag 00GQ)}.
Thus $|U| \to |M|$ is surjective.

\medskip\noindent
Let $k$ be an algebraically closed field and let $C \to k$ be a ring map.
Since surjective morphisms are preserved under base change
(Morphisms, Lemma \href{https://stacks.math.columbia.edu/tag/01S1}{lemma base change surjective (Tag 01S1)})
we see that $A \otimes_C k$ is not zero. Now $k \subset A \otimes_C k$ is a
nonzero integral extension. Hence any residue field of $A \otimes_C k$
is an algebraic extension of $k$, hence equal to $k$. Thus we see that
$U(k) \to M(k)$ is surjective.

\medskip\noindent
Let $a_0, a_1 : A \to k$ be two ring maps. If there exists a ring map
$b : B \to k$ such that $a_0 = b \circ t^\sharp$ and $a_1 = b \circ s^\sharp$
then we see that $a_0|_C = a_1|_C$ by definition. Thus the map
$U(k) \to M(k)$ equalizes the two maps $R(k) \to U(k)$.
Conversely, suppose that $a_0|_C = a_1|_C$. Let us name this algebra
map $c : C \to k$. Consider the diagram
$$
\xymatrix{
& &
B \ar@{-->}[lld] \\
k & &
A
\ar@<0.5ex>[ll]^{a_0}
\ar@<-0.5ex>[ll]_{a_1}
\ar@<1ex>[u]
\ar@<-1ex>[u] \\
& &
C \ar[u] \ar[llu]^c
}
$$
If we can construct a dotted arrow making the diagram commute, then
the proof of part (2) of the lemma is complete. Since $s : A \to B$ is finite
there exist finitely many ring maps
$b_1, \ldots, b_n : B \to k$ such that $b_i \circ s^\sharp = a_1$.
If the dotted arrow does not exist, then we see that none of the
$a'_i = b_i \circ t^\sharp$, $i = 1, \ldots, n$ is equal to $a_0$.
Hence the maximal ideals
$$
\mathfrak m'_i = \Ker(a_i' \otimes 1 : A \otimes_C k \to k)
$$
of $A \otimes_C k$ are distinct from
$\mathfrak m = \Ker(a_0 \otimes 1 : A \otimes_C k \to k)$.
By Algebra, Lemma \href{https://stacks.math.columbia.edu/tag/00DS}{lemma silly (Tag 00DS)} we would get an element
$f \in A \otimes_C k$ with $f \in \mathfrak m$, but
$f \not \in \mathfrak m_i'$ for $i = 1, \ldots, n$.
Consider the norm
$$
g = \text{Norm}_{s^\sharp \otimes 1}(t^\sharp \otimes 1(f))
\in
A \otimes_C k
$$
By Lemma \ref{groupoids-lemma-determinant-trick} this lies in the invariants
$C^1 \subset A \otimes_C k$ of the base change
groupoid (base change via the map $c : C \to k$). On the one hand,
$a_1(g) \in k^*$ since
the value of $t^\sharp(f)$ at all the points (which correspond to
$b_1, \ldots, b_n$) lying over $a_1$ is
invertible (insert future reference on property determinant here).
On the other hand, since $f \in \mathfrak m$, we see that
$f$ is not a unit, hence $t^\sharp(f)$ is not a unit
(as $t^\sharp \otimes 1$ is faithfully flat),
hence its norm is not a unit (insert future reference
on property determinant here). We conclude that $C^1$ contains
an element which is not nilpotent
and not a unit. We will now show that this leads to a contradiction.
Namely, apply Lemma \ref{groupoids-lemma-invariants-base-change}
to the map $c : C \to C' = k$, then
we see that the map of $k$ into the invariants $C^1$ is injective
and moreover, that for any element $x \in C^1$ there exists an integer
$n > 0$ such that $x^n \in k$. Hence every element of $C^1$ is
either a unit or nilpotent.

\medskip\noindent
We still have to finish the proof of (1). We already know that
$|U| \to |M|$ is surjective. It is clear that $|U| \to |M|$ is
$|R|$-invariant. Finally, suppose $u_0, u_1 \in U$ maps to the same
point $m \in M$. Then the induced field extensions $\kappa(u_0)/\kappa(m)$
and $\kappa(u_1)/\kappa(m)$ are algebraic (as $A$ is integral over $C$
as used above). Hence if $k$ is an algebraic closure of $\kappa(m)$
then we can find $\kappa(m)$-embeddings $\overline{u}_0 : \kappa(u_0) \to k$
and $\overline{u}_1 : \kappa(u_1) \to k$. These determine $k$-valued
points $\overline{u}_0, \overline{u}_1 \in U(k)$ mapping to the same
point of $M(k)$. By part (2) we see that there exists a point
$\overline{r} \in R(k)$ with $s(\overline{r}) = \overline{u}_0$ and
$t(\overline{r}) = \overline{u}_1$. The image $r \in R$ of $\overline{r}$
is a point with $s(r) = u_0$ and $t(r) = u_1$ as desired.
\end{proof}


\begin{lemma}
\label{groupoids-lemma-basis}
Let $S$ be a scheme.
Let $(U, R, s, t, c)$ be a groupoid scheme over $S$.
Assume
\begin{enumerate}
\item $U = \Spec(A)$, and $R = \Spec(B)$ are affine, and
\item there exist elements $x_i \in A$, $i \in I$ such that
$B = \bigoplus_{i \in I} s^\sharp(A)t^\sharp(x_i)$.
\end{enumerate}
Then $A = \bigoplus_{i\in I} Cx_i$, and $B \cong A \otimes_C A$
where $C \subset A$ is the $R$-invariant
functions on $U$ as in (\ref{groupoids-equation-invariants}).
\end{lemma}

\begin{proof}
During this proof we will write $s, t : A \to B$ instead of
$s^\sharp, t^\sharp$, and similarly $c : B \to B \otimes_{s, A, t} B$.
We write $p_0 : B \to B \otimes_{s, A, t} B$, $b \mapsto b \otimes 1$ and
$p_1 : B \to B \otimes_{s, A, t} B$, $b \mapsto 1 \otimes b$. By
Lemma \ref{groupoids-lemma-diagram-pull}
and the definition of $C$ we have the following
commutative diagram
$$
\xymatrix{
B \otimes_{s, A, t} B &
B \ar@<-1ex>[l]_-c \ar@<1ex>[l]^-{p_0} &
A \ar[l]^t \\
B \ar[u]^{p_1} &
A \ar@<-1ex>[l]_s \ar@<1ex>[l]^t \ar[u]_s &
C \ar[u] \ar[l]
}
$$
Moreover the tow left squares are cocartesian in the category of rings, and
the top row is isomorphic to the diagram
$$
\xymatrix{
B \otimes_{t, A, t} B &
B \ar@<-1ex>[l]_-{p_1} \ar@<1ex>[l]^-{p_0} &
A \ar[l]^t
}
$$
which is an equalizer diagram according to
Descent, Lemma \href{https://stacks.math.columbia.edu/tag/023M}{lemma ff exact (Tag 023M)} because condition (2) implies
in particular that $s$ (and hence also then isomorphic arrow $t$)
is faithfully flat.
The lower row is an equalizer diagram by definition of $C$.
We can use the $x_i$ and get a commutative diagram
$$
\xymatrix{
B \otimes_{s, A, t} B &
B \ar@<-1ex>[l]_-c \ar@<1ex>[l]^-{p_0} &
A \ar[l]^t \\
\bigoplus_{i \in I} B x_i \ar[u]^{p_1} &
\bigoplus_{i \in I} A x_i \ar@<-1ex>[l]_s \ar@<1ex>[l]^t \ar[u]_s &
\bigoplus_{i \in I} C x_i \ar[u] \ar[l]
}
$$
where in the right vertical arrow we map $x_i$ to $x_i$,
in the middle vertical arrow we map $x_i$ to $t(x_i)$ and
in the left vertical arrow we map $x_i$ to
$c(t(x_i)) = t(x_i) \otimes 1 = p_0(t(x_i))$ (equality by the commutativity
of the top part of the diagram in Lemma \ref{groupoids-lemma-diagram}). Then the diagram
commutes. Moreover the middle vertical arrow is an isomorphism
by assumption. Since the left two squares are cocartesian we
conclude that also the left vertical arrow is an isomorphism.
On the other hand, the horizontal rows are exact (i.e., they are
equalizers). Hence we conclude that also the right vertical arrow
is an isomorphism.
\end{proof}


\begin{lemma}
\label{groupoids-lemma-diagram-pull}
Let $S$ be a scheme.
Let $(U, R, s, t, c, e, i)$ be a groupoid over $S$.
The diagram
\begin{equation}
\label{groupoids-equation-pull}
\xymatrix{
R \times_{t, U, t} R
\ar@<1ex>[r]^-{\text{pr}_1} \ar@<-1ex>[r]_-{\text{pr}_0}
\ar[d]_{(\text{pr}_0, c \circ (i, 1))} &
R \ar[r]^t \ar[d]^{\text{id}_R} &
U \ar[d]^{\text{id}_U} \\
R \times_{s, U, t} R
\ar@<1ex>[r]^-c \ar@<-1ex>[r]_-{\text{pr}_0} \ar[d]_{\text{pr}_1} &
R \ar[r]^t \ar[d]^s &
U \\
R \ar@<1ex>[r]^s \ar@<-1ex>[r]_t &
U
}
\end{equation}
is commutative. The two top rows are isomorphic via the vertical maps given.
The two lower left squares are cartesian.
\end{lemma}

\begin{proof}
The commutativity of the diagram follows from the axioms of a groupoid.
Note that, in terms of groupoids, the top left vertical arrow assigns to
a pair of morphisms $(\alpha, \beta)$ with the same target, the pair
of morphisms $(\alpha, \alpha^{-1} \circ \beta)$. In any groupoid
this defines a bijection between
$\text{Arrows} \times_{t, \text{Ob}, t} \text{Arrows}$
and
$\text{Arrows} \times_{s, \text{Ob}, t} \text{Arrows}$. Hence the second
assertion of the lemma.
The last assertion follows from Lemma \ref{groupoids-lemma-diagram}.
\end{proof}


\begin{lemma}
\label{groupoids-lemma-diagram}
Let $S$ be a scheme.
Let $(U, R, s, t, c)$ be a groupoid over $S$.
In the commutative diagram
$$
\xymatrix{
& U & \\
R \ar[d]_s \ar[ru]^t &
R \times_{s, U, t} R
\ar[l]^-{\text{pr}_0} \ar[d]^{\text{pr}_1} \ar[r]_-c &
R \ar[d]^s \ar[lu]_t \\
U & R \ar[l]_t \ar[r]^s & U
}
$$
the two lower squares are fibre product squares.
Moreover, the triangle on top (which is really a square)
is also cartesian.
\end{lemma}

\begin{proof}
Omitted.
Exercise in the definitions and the functorial point of
view in algebraic geometry.
\end{proof}


\begin{lemma}
\label{groupoids-lemma-criterion-quotient-representable}
In the situation of Definition \ref{groupoids-definition-quotient-sheaf}.
Assume there is a scheme $M$, and a morphism $U \to M$ such that
\begin{enumerate}
\item the morphism $U \to M$ equalizes $s, t$,
\item the morphism $U \to M$ induces a surjection of sheaves
$h_U \to h_M$ in the $\tau$-topology, and
\item the induced map $(t, s) : R \to U \times_M U$ induces a
surjection of sheaves $h_R \to h_{U \times_M U}$ in the $\tau$-topology.
\end{enumerate}
In this case $M$ represents the quotient sheaf $U/R$.
\end{lemma}

\begin{proof}
Condition (1) says that $h_U \to h_M$ factors through $U/R$.
Condition (2) says that $U/R \to h_M$ is surjective as a map of sheaves.
Condition (3) says that $U/R \to h_M$ is injective as a map of sheaves.
Hence the lemma follows.
\end{proof}


\begin{lemma}
\label{algebra-lemma-semi-local-module-basis-in-submodule}
Let $R$ be a local ring with maximal ideal $\mathfrak m$ and
infinite residue field.
Let $R \to S$ be a ring map.
Let $M$ be an $S$-module and let $N \subset M$ be an $R$-submodule.
Assume
\begin{enumerate}
\item $S$ is semi-local and $\mathfrak mS$ is contained in the
Jacobson radical of $S$,
\item $M$ is a finite free $S$-module, and
\item $N$ generates $M$ as an $S$-module.
\end{enumerate}
Then $N$ contains an $S$-basis of $M$.
\end{lemma}

\begin{proof}
Assume $M$ is free of rank $n$. Let $I \subset S$ be the Jacobson radical.
By Nakayama's Lemma \href{https://stacks.math.columbia.edu/tag/00DV}{lemma NAK (Tag 00DV)} a sequence of elements
$m_1, \ldots, m_n$ is a basis for $M$ if and only if
$\overline{m}_i \in M/IM$ generate $M/IM$. Hence we may replace
$M$ by $M/IM$, $N$ by $N/(N \cap IM)$, $R$ by $R/\mathfrak m$,
and $S$ by $S/IS$. In this case we see that $S$ is a finite product
of fields $S = k_1 \times \ldots \times k_r$ and
$M = k_1^{\oplus n} \times \ldots \times k_r^{\oplus n}$.
The fact that $N \subset M$ generates $M$ as an $S$-module
means that there exist $x_j \in N$ such that a linear combination
$\sum a_j x_j$ with $a_j \in S$ has a nonzero component in each
factor $k_i^{\oplus n}$.
Because $R = k$ is an infinite field, this means that also
some linear combination $y = \sum c_j x_j$ with $c_j \in k$ has a
nonzero component in each factor. Hence $y \in N$ generates a
free direct summand $Sy \subset M$. By induction on $n$ the result
holds for $M/Sy$ and the submodule $\overline{N} = N/(N \cap Sy)$.
In other words there exist $\overline{y}_2, \ldots, \overline{y}_n$
in $\overline{N}$ which (freely) generate $M/Sy$. Then
$y, y_2, \ldots, y_n$ (freely) generate $M$ and we win.
\end{proof}


\begin{lemma}
\label{algebra-lemma-flat-local-given-residue-field}
Let $(R, \mathfrak m, k)$ be a local ring. Let $K/k$ be a field
extension. There exists a local ring $(R', \mathfrak m', k')$, a flat local
ring map $R \to R'$ such that $\mathfrak m' = \mathfrak mR'$ and such that
$k'$ is isomorphic to $K$ as an extension of $k$.
\end{lemma}

\begin{proof}
Suppose that $k' = k(\alpha)$ is a monogenic extension of $k$.
Then $k'$ is the residue field of a flat local extension $R \subset R'$
as in the lemma. Namely, if $\alpha$ is transcendental over $k$, then we let
$R'$ be the localization of $R[x]$ at the prime $\mathfrak mR[x]$.
If $\alpha$ is algebraic with minimal polynomial
$T^d + \sum \overline{\lambda}_iT^{d - i}$, then we let
$R' = R[T]/(T^d + \sum \lambda_i T^{d - i})$.

\medskip\noindent
Consider the collection of triples $(k', R \to R', \phi)$, where
$k \subset k' \subset K$ is a subfield,
$R \to R'$ is a local ring map as in the lemma, and
$\phi : R' \to k'$ induces an isomorphism $R'/\mathfrak mR' \cong k'$
of $k$-extensions. These form a ``big'' category $\mathcal{C}$ with morphisms
$(k_1, R_1, \phi_1) \to (k_2, R_2, \phi_2)$
given by ring maps $\psi : R_1 \to R_2$ such that
$$
\xymatrix{
R_1 \ar[d]_\psi \ar[r]_{\phi_1} & k_1 \ar[r] & K \ar@{=}[d] \\
R_2 \ar[r]^{\phi_2} & k_2 \ar[r] & K
}
$$
commutes. This implies that $k_1 \subset k_2$.

\medskip\noindent
Suppose that $I$ is a directed set, and
$((R_i, k_i, \phi_i), \psi_{ii'})$ is a system over $I$, see
Categories, Section \href{https://stacks.math.columbia.edu/tag/002Z}{section posets limits (Tag 002Z)}.
In this case we can consider
$$
R' = \colim_{i \in I} R_i
$$
This is a local ring with maximal ideal $\mathfrak mR'$, and
residue field $k' = \bigcup_{i \in I} k_i$. Moreover, the ring
map $R \to R'$ is flat as it is a colimit of flat maps (and tensor
products commute with directed colimits).
Hence we see that $(R', k', \phi')$ is an ``upper bound'' for the system.

\medskip\noindent
An almost trivial application of Zorn's Lemma would finish the proof
if $\mathcal{C}$ was a set, but it isn't.
(Actually, you can make this work by finding a reasonable bound on the
cardinals of the local rings occurring.)
To get around this problem we choose a well ordering on $K$.
For $x \in K$ we let $K(x)$ be the subfield of $K$ generated
by all elements of $K$ which are $\leq x$.
By transfinite recursion on $x \in K$ we will produce ring maps
$R \subset R(x)$ as in the lemma with residue field extension
$K(x)/k$. Moreover, by construction we will have that
$R(x)$ will contain $R(y)$ for all $y \leq x$.
Namely, if $x$ has a predecessor $x'$, then $K(x) = K(x')[x]$
and hence we can let $R(x') \subset R(x)$ be the local ring extension
constructed in the first paragraph of the proof. If $x$ does not
have a predecessor, then we first set
$R'(x) = \colim_{x' < x} R(x')$ as in the third paragraph
of the proof. The residue field of $R'(x)$ is $K'(x) = \bigcup_{x' < x} K(x')$.
Since $K(x) = K'(x)[x]$ we see that we can use the construction of the
first paragraph of the proof to produce $R'(x) \subset R(x)$.
This finishes the proof of the lemma.
\end{proof}


Let $\tau \in \{Zariski, \etale, fppf, smooth, syntomic\}$.
Let $S$ be a scheme.
Let $j : R \to U \times_S U$ be a pre-relation over $S$.
Say $U, R, S$ are objects of a $\tau$-site $\Sch_\tau$
(see Topologies, Section \href{https://stacks.math.columbia.edu/tag/020M}{section procedure (Tag 020M)}).
Then we can consider the functors
$$
h_U, h_R :
(\Sch/S)_\tau^{opp}
\longrightarrow
\textit{Sets}.
$$
These are sheaves, see
Descent, Lemma \href{https://stacks.math.columbia.edu/tag/023Q}{lemma fpqc universal effective epimorphisms (Tag 023Q)}.
The morphism $j$ induces a map $j : h_R \to h_U \times h_U$.
For each object $T \in \Ob((\Sch/S)_\tau)$
we can take the equivalence relation $\sim_T$ generated by
$j(T) : R(T) \to U(T) \times U(T)$ and consider the quotient.
Hence we get a presheaf
\begin{equation}
\label{groupoids-equation-quotient-presheaf}
(\Sch/S)_\tau^{opp}
\longrightarrow
\textit{Sets}, \quad
T \longmapsto U(T)/\sim_T
\end{equation}


\noindent
Recall that a groupoid is a category in which every morphism
is an isomorphism, see
Categories, Definition \ref{categories-definition-groupoid}.
Hence a groupoid has a set of objects $\text{Ob}$,
a set of arrows $\text{Arrows}$, a {\it source} and {\it target}
map $s, t : \text{Arrows} \to \text{Ob}$, and a {\it composition law}
$c : \text{Arrows} \times_{s, \text{Ob}, t} \text{Arrows}
\to \text{Arrows}$.
These maps satisfy exactly the following axioms
\begin{enumerate}
\item (associativity) $c \circ (1, c) = c \circ (c, 1)$ as maps
$\text{Arrows} \times_{s, \text{Ob}, t}
\text{Arrows} \times_{s, \text{Ob}, t}
\text{Arrows} \to \text{Arrows}$,
\item (identity) there exists a map $e : \text{Ob} \to \text{Arrows}$
such that
\begin{enumerate}
\item $s \circ e = t \circ e = \text{id}$ as maps $\text{Ob} \to \text{Ob}$,
\item $c \circ (1, e \circ s) = c \circ (e \circ t, 1) = 1$
as maps $\text{Arrows} \to \text{Arrows}$,
\end{enumerate}
\item (inverse) there exists a map $i : \text{Arrows} \to \text{Arrows}$
such that
\begin{enumerate}
\item $s \circ i = t$, $t \circ i = s$ as maps $\text{Arrows} \to \text{Ob}$,
and
\item $c \circ (1, i) = e \circ t$ and $c \circ (i, 1) = e \circ s$
as maps $\text{Arrows} \to \text{Arrows}$.
\end{enumerate}
\end{enumerate}
If this is the case the maps $e$ and $i$ are uniquely determined and
$i$ is a bijection. Note that if $(\text{Ob}', \text{Arrows}', s', t', c')$
is a second groupoid category, then a functor
$f : (\text{Ob}, \text{Arrows}, s, t, c) \to
(\text{Ob}', \text{Arrows}', s', t', c')$
is given by a pair of set maps $f : \text{Ob} \to \text{Ob}'$ and
$f : \text{Arrows} \to \text{Arrows}'$ such that
$s' \circ f = f \circ s$, $t' \circ f = f \circ t$, and
$c' \circ (f, f) = f \circ c$. The compatibility with identity and
inverse is automatic. We will use this below.
(Warning: The compatibility with identity
has to be imposed in the case of general categories.)


\begin{lemma}
\label{lemma-descend-equality}
Notation and assumptions as in Situation \ref{situation-descent}.
Let $f_0 : Y_0 \to Z_0$ be a morphism of algebraic spaces over $X_0$.
Assume (a) $Y_0 \to X_0$ and $Z_0 \to X_0$ are representable, (b)
$Y_0$, $Z_0$ quasi-compact and quasi-separated, (c)
$f_0$ locally of finite presentation, and
(d) $Y_0 \times_{X_0} X \to Z_0 \times_{X_0} X$ an isomorphism.
Then there exists an $i \geq 0$ such that
$Y_0 \times_{X_0} X_i \to Z_0 \times_{X_0} X_i$ is an isomorphism.
\end{lemma}

\begin{proof}
Choose an affine scheme $U_0$ and a surjective \'etale morphism $U_0 \to X_0$.
Set $U_i = U_0 \times_{X_0} X_i$ and $U = U_0 \times_{X_0} X$.
Apply Limits, Lemma \ref{limits-lemma-descend-isomorphism}
to see that $Y_0 \times_{X_0} U_i \to Z_0 \times_{X_0} U_i$
is an isomorphism of schemes for some $i \geq 0$ (details omitted).
As $U_i \to X_i$ is surjective \'etale, it follows that
$Y_0 \times_{X_0} X_i \to Z_0 \times_{X_0} X_i$ is an isomorphism
(details omitted).
\end{proof}


\begin{lemma}
\label{limits-lemma-descend-isomorphism}
Notation and assumptions as in Situation \ref{limits-situation-descent-property}.
If
\begin{enumerate}
\item $f$ is an isomorphism, and
\item $f_0$ is locally of finite presentation,
\end{enumerate}
then $f_i$ is an isomorphism for some $i \geq 0$.
\end{lemma}

\begin{proof}
By Lemmas \ref{limits-lemma-descend-etale} and
\ref{limits-lemma-descend-closed-immersion-finite-presentation}
we can find an $i$ such that $f_i$ is flat and a closed immersion.
Then $f_i$ identifies $X_i$ with an open and closed subscheme of
$Y_i$, see Morphisms, Lemma
\href{https://stacks.math.columbia.edu/tag/0819}{lemma flat closed immersions finite presentation (Tag 0819)}.
By assumption the image of $Y \to Y_i$ maps into $f_i(X_i)$.
Thus by Lemma \href{https://stacks.math.columbia.edu/tag/05F4}{lemma limit contained in constructible (Tag 05F4)}
we find that $Y_{i'}$ maps into $f_i(X_i)$ for some $i' \geq i$.
It follows that $X_{i'} \to Y_{i'}$ is surjective and we win.
\end{proof}


\begin{lemma}
\label{limits-lemma-descend-affine-finite-presentation}
Notation and assumptions as in Situation \ref{limits-situation-descent-property}.
If $f$ is affine, then there exists an index $i \geq 0$
such that $f_i$ is affine.
\end{lemma}

\begin{proof}
Let $Y_0 = \bigcup_{j = 1, \ldots, m} V_{j, 0}$ be a finite affine
open covering. Set $U_{j, 0} = f_0^{-1}(V_{j, 0})$. For $i \geq 0$
we denote $V_{j, i}$ the inverse image of $V_{j, 0}$ in $Y_i$ and
$U_{j, i} = f_i^{-1}(V_{j, i})$. Similarly we have
$U_j = f^{-1}(V_j)$. Then $U_j = \lim_{i \geq 0} U_{j, i}$
(see Lemma \ref{limits-lemma-directed-inverse-system-has-limit}).
Since $U_j$ is affine by assumption we see that
each $U_{j, i}$ is affine for $i$ large enough, see
Lemma \ref{limits-lemma-limit-affine}. As there are finitely many $j$ we
can pick an $i$ which works for all $j$. Thus $f_i$ is
affine for $i$ large enough, see
Morphisms, Lemma \href{https://stacks.math.columbia.edu/tag/01S8}{lemma characterize affine (Tag 01S8)}.
\end{proof}


\begin{lemma}
\label{limits-lemma-descend-closed-immersion-finite-presentation}
Notation and assumptions as in Situation \ref{limits-situation-descent-property}.
If
\begin{enumerate}
\item $f$ is a closed immersion, and
\item $f_0$ is locally of finite type,
\end{enumerate}
then there exists an $i \geq 0$ such that $f_i$ is a closed immersion.
\end{lemma}

\begin{proof}
A closed immersion is affine, see
Morphisms, Lemma \href{https://stacks.math.columbia.edu/tag/01SE}{lemma closed immersion affine (Tag 01SE)}.
Hence by Lemma \ref{limits-lemma-descend-affine-finite-presentation} above
after increasing $0$ we may assume that $f_0$ is affine.
By writing $Y_0$ as a finite union of affines we reduce to proving
the result when $X_0$ and $Y_0$ are affine and map
into a common affine $W \subset S_0$. The corresponding algebra
statement is a consequence of
Algebra, Lemma \href{https://stacks.math.columbia.edu/tag/07RH}{lemma colimit surjective (Tag 07RH)}.
\end{proof}


\begin{lemma}
\label{limits-lemma-descend-etale}
Notation and assumptions as in Situation \ref{limits-situation-descent-property}.
If
\begin{enumerate}
\item $f$ is \'etale,
\item $f_0$ is locally of finite presentation,
\end{enumerate}
then $f_i$ is \'etale for some $i \geq 0$.
\end{lemma}

\begin{proof}
Being \'etale is local on the source and the target (Morphisms,
Lemma \href{https://stacks.math.columbia.edu/tag/02GJ}{lemma etale characterize (Tag 02GJ)}) hence we may assume
$S_0, X_0, Y_0$ affine (details omitted). The corresponding algebra fact is
Algebra, Lemma \ref{algebra-lemma-colimit-etale}.
\end{proof}


\begin{lemma}
\label{algebra-lemma-colimit-etale}
Let $A = \colim_{i \in I} A_i$ be a directed colimit of rings.
Let $0 \in I$ and $\varphi_0 : B_0 \to C_0$ a map of $A_0$-algebras.
Assume
\begin{enumerate}
\item $A \otimes_{A_0} B_0 \to A \otimes_{A_0} C_0$ is \'etale,
\item $B_0 \to C_0$ is of finite presentation.
\end{enumerate}
Then for some $i \geq 0$ the map
$A_i \otimes_{A_0} B_0 \to A_i \otimes_{A_0} C_0$
is \'etale.
\end{lemma}

\begin{proof}
Write $C_0 = B_0[x_1, \ldots, x_n]/(f_{1, 0}, \ldots, f_{m, 0})$.
Write $B_i = A_i \otimes_{A_0} B_0$ and $C_i = A_i \otimes_{A_0} C_0$.
Note that $C_i = B_i[x_1, \ldots, x_n]/(f_{1, i}, \ldots, f_{m, i})$
where $f_{j, i}$ is the image of $f_{j, 0}$ in the polynomial ring
over $B_i$. Write $B = A \otimes_{A_0} B_0$ and $C = A \otimes_{A_0} C_0$.
Note that $C = B[x_1, \ldots, x_n]/(f_1, \ldots, f_m)$
where $f_j$ is the image of $f_{j, 0}$ in the polynomial ring
over $B$. The assumption is that the map
$$
\text{d} :
(f_1, \ldots, f_m)/(f_1, \ldots, f_m)^2
\longrightarrow
\bigoplus C \text{d}x_k
$$
is an isomorphism. Thus for sufficiently large $i$ we can find elements
$$
\xi_{k, i} \in (f_{1, i}, \ldots, f_{m, i})/(f_{1, i}, \ldots, f_{m, i})^2
$$
with $\text{d}\xi_{k, i} = \text{d}x_k$ in $\bigoplus C_i \text{d}x_k$.
Moreover, on increasing $i$ if necessary, we see that
$\sum (\partial f_{j, i}/\partial x_k) \xi_{k, i} =
f_{j, i} \bmod (f_{1, i}, \ldots, f_{m, i})^2$
since this is true in the limit. Then this $i$ works.
\end{proof}


\begin{lemma}
\label{limits-lemma-descend-monomorphism}
Notation and assumptions as in Situation \ref{limits-situation-descent-property}.
If
\begin{enumerate}
\item $f$ is a monomorphism, and
\item $f_0$ is locally of finite type,
\end{enumerate}
then $f_i$ is a monomorphism for some $i \geq 0$.
\end{lemma}

\begin{proof}
Recall that a morphism of schemes $V \to W$ is a monomorphism if and
only if the diagonal $V \to V \times_W V$ is an isomorphism
(Schemes, Lemma \href{https://stacks.math.columbia.edu/tag/01L3}{lemma monomorphism (Tag 01L3)}).
The morphism $X_0 \to X_0 \times_{Y_0} X_0$ is locally of finite
presentation by
Morphisms, Lemma \href{https://stacks.math.columbia.edu/tag/0818}{lemma diagonal morphism finite type (Tag 0818)}.
Since $X_0 \times_{Y_0} X_0$ is quasi-compact and quasi-separated
(Schemes, Remark \href{https://stacks.math.columbia.edu/tag/0816}{remark quasi compact and quasi separated (Tag 0816)})
we conclude from
Lemma \ref{limits-lemma-descend-isomorphism}
that $\Delta_i : X_i \to X_i \times_{Y_i} X_i$ is an isomorphism for
some $i \geq 0$. For this $i$ the morphism $f_i$ is a monomorphism.
\end{proof}


\begin{lemma}
\label{limits-lemma-approximate}
Let $S$ be a quasi-compact and quasi-separated scheme. Let $V \subset S$
be a quasi-compact open. Let $I$ be a directed set
and let $(V_i, f_{ii'})$ be an inverse system of schemes over $I$
with affine transition maps, with each $V_i$ of finite type over $\mathbf{Z}$,
and with $V = \lim V_i$. Then there exist
\begin{enumerate}
\item a directed set $J$,
\item an inverse system of schemes $(S_j, g_{jj'})$ over $J$,
\item an order preserving map $\alpha : J \to I$,
\item open subschemes $V'_j \subset S_j$, and
\item isomorphisms $V'_j \to V_{\alpha(j)}$
\end{enumerate}
such that
\begin{enumerate}
\item the transition morphisms $g_{jj'} : S_j \to S_{j'}$ are affine,
\item each $S_j$ is of finite type over $\mathbf{Z}$,
\item $g_{jj'}^{-1}(V'_{j'}) = V'_j$,
\item $S = \lim S_j$ and $V = \lim V'_j$, and
\item the diagrams
$$
\vcenter{
\xymatrix{
V \ar[d] \ar[rd] \\
V'_j \ar[r] & V_{\alpha(j)}
}
}
\quad\text{and}\quad
\vcenter{
\xymatrix{
V'_j \ar[r] \ar[d] & V_{\alpha(j)} \ar[d] \\
V'_{j'} \ar[r] & V_{\alpha(j')}
}
}
$$
are commutative.
\end{enumerate}
\end{lemma}

\begin{proof}
Set $Z = S \setminus V$. Choose affine opens $U_1, \ldots, U_m \subset S$
such that $Z \subset \bigcup_{l = 1, \ldots, m} U_l$. Consider the opens
$$
V \subset V \cup U_1 \subset V \cup U_1 \cup U_2 \subset
\ldots \subset V \cup \bigcup\nolimits_{l = 1, \ldots, m} U_l = S
$$
If we can prove the lemma successively for each of the cases
$$
V \cup U_1 \cup \ldots \cup U_l
\subset
V \cup U_1 \cup \ldots \cup U_{l + 1}
$$
then the lemma will follow for $V \subset S$. In each case we are adding
one affine open. Thus we may assume
\begin{enumerate}
\item $S = U \cup V$,
\item $U$ affine open in $S$,
\item $V$ quasi-compact open in $S$, and
\item $V = \lim_i V_i$ with $(V_i, f_{ii'})$
an inverse system over a directed set $I$, each $f_{ii'}$
affine and each $V_i$ of finite type over $\mathbf{Z}$.
\end{enumerate}
Denote $f_i : V \to V_i$ the projections.
Set $W = U \cap V$. As $S$ is quasi-separated, this is a quasi-compact open
of $V$. By Lemma \ref{limits-lemma-descend-opens}
(and after shrinking $I$) we may assume that there exist
opens $W_i \subset V_i$ such that $f_{ii'}^{-1}(W_{i'}) = W_i$
and such that $f_i^{-1}(W_i) = W$. Since $W$ is a quasi-compact open
of $U$ it is quasi-affine. Hence we may assume (after shrinking $I$ again)
that $W_i$ is quasi-affine for all $i$, see
Lemma \ref{limits-lemma-limit-quasi-affine}.

\medskip\noindent
Write $U = \Spec(B)$. Set $R = \Gamma(W, \mathcal{O}_W)$,
and $R_i = \Gamma(W_i, \mathcal{O}_{W_i})$.
By Lemma \ref{limits-lemma-descend-section} we have $R = \colim_i R_i$.
Now we have the maps of rings
$$
\xymatrix{
B \ar[r]_s & R \\
& R_i \ar[u]_{t_i}
}
$$
We set $B_i = \{(b, r) \in B \times R_i \mid s(b) = t_i(r)\}$ so that we
have a cartesian diagram
$$
\xymatrix{
B \ar[r]_s & R \\
B_i \ar[u] \ar[r] & R_i \ar[u]_{t_i}
}
$$
for each $i$. The transition maps $R_i \to R_{i'}$ induce maps
$B_i \to B_{i'}$. It is clear that $B = \colim_i B_i$.
In the next paragraph we show that for all sufficiently large $i$
the composition $W_i \to \Spec(R_i) \to \Spec(B_i)$ is an open immersion.

\medskip\noindent
As $W$ is a quasi-compact open of $U = \Spec(B)$
we can find a finitely many elements $g_l \in B$, $l = 1, \ldots, m$
such that $D(g_l) \subset W$ and such that
$W = \bigcup_{l = 1, \ldots, m} D(g_l)$.
Note that this implies $D(g_l) = W_{s(g_l)}$ as open subsets of $U$,
where $W_{s(g_l)}$ denotes the largest open subset of $W$ on which
$s(g_l)$ is invertible. Hence
$$
B_{g_l} =
\Gamma(D(g_l), \mathcal{O}_U) =
\Gamma(W_{s(g_l)}, \mathcal{O}_W) = R_{s(g_l)},
$$
where the last equality is
Properties, Lemma \href{https://stacks.math.columbia.edu/tag/01P7}{lemma invert f sections (Tag 01P7)}.
Since $W_{s(g_l)}$ is affine this also
implies that $D(s(g_l)) = W_{s(g_l)}$ as open subsets of $\Spec(R)$.
Since $R = \colim_i R_i$ we can (after shrinking $I$)
assume there exist $g_{l, i} \in R_i$ for all $i \in I$ such that
$s(g_l) = t_i(g_{l, i})$. Of course we choose the $g_{l, i}$
such that $g_{l, i}$ maps to $g_{l, i'}$ under the transition maps
$R_i \to R_{i'}$. Then, by Lemma \ref{limits-lemma-descend-opens} we can
(after shrinking $I$ again)
assume the corresponding opens $D(g_{l, i}) \subset \Spec(R_i)$
are contained in $W_i$ for $l = 1, \ldots, m$ and cover $W_i$.
We conclude that the morphism $W_i \to \Spec(R_i) \to \Spec(B_i)$
is an open immersion, see Lemma \ref{limits-lemma-diagram}.

\medskip\noindent
By Lemma \ref{limits-lemma-quasi-affine-finite-type-over-Z}
we can write $B_i$ as a directed colimit of subalgebras
$A_{i, p} \subset B_i$, $p \in P_i$ each
of finite type over $\mathbf{Z}$ and such that $W_i$ is
identified with an open subscheme of $\Spec(A_{i, p})$.
Let $S_{i, p}$ be the scheme obtained by glueing
$V_i$ and $\Spec(A_{i, p})$ along the open $W_i$, see
Schemes, Section \href{https://stacks.math.columbia.edu/tag/01JA}{section glueing schemes (Tag 01JA)}.
Here is the resulting commutative diagram of schemes:
$$
\xymatrix{
& & V \ar[lld] \ar[d] & W \ar[l] \ar[lld] \ar[d] \\
V_i \ar[d] & W_i \ar[l] \ar[d] & S \ar[lld] & U \ar[lld] \ar[l] \\
S_{i, p} & \Spec(A_{i, p}) \ar[l]
}
$$
The morphism $S \to S_{i, p}$ arises because the upper right
square is a pushout in the category of schemes.
Note that $S_{i, p}$ is of finite type over $\mathbf{Z}$ since
it has a finite affine open covering whose members are
spectra of finite type $\mathbf{Z}$-algebras.
We define a preorder on $J = \coprod_{i \in I} P_i$
by the rule $(i', p') \geq (i, p)$ if and only if
$i' \geq i$ and the map $B_i \to B_{i'}$ maps $A_{i, p}$ into
$A_{i', p'}$. This is exactly the condition needed to
define a morphism $S_{i', p'} \to S_{i, p}$: namely make a commutative
diagram as above using the transition morphisms $V_{i'} \to V_i$
and $W_{i'} \to W_i$ and
the morphism $\Spec(A_{i', p'}) \to \Spec(A_{i, p})$ induced
by the ring map $A_{i, p} \to A_{i', p'}$. The relevant commutativities
have been built into the constructions.
We claim that $S$ is the directed limit of the schemes $S_{i, p}$.
Since by construction the schemes $V_i$ have limit $V$ this boils
down to the fact that $B$ is the limit of the rings $A_{i, p}$
which is true by construction. The map $\alpha : J \to I$ is given
by the rule $j = (i, p) \mapsto i$. The open subscheme $V'_j$ is
just the image of $V_i \to S_{i, p}$ above. The commutativity of
the diagrams in (5) is clear from the construction.
This finishes the proof of the lemma.
\end{proof}


\begin{lemma}
\label{limits-lemma-descend-opens}
In Situation \ref{limits-situation-descent} we have the following:
\begin{enumerate}
\item Given any quasi-compact open $V \subset S = \lim_i S_i$
there exists an $i \in I$ and a quasi-compact open $V_i \subset S_i$
such that $f_i^{-1}(V_i) = V$.
\item Given $V_i \subset S_i$ and $V_{i'} \subset S_{i'}$
quasi-compact opens such that $f_i^{-1}(V_i) = f_{i'}^{-1}(V_{i'})$
there exists an index $i'' \geq i, i'$ such that
$f_{i''i}^{-1}(V_i) = f_{i''i'}^{-1}(V_{i'})$.
\item If $V_{1, i}, \ldots, V_{n, i} \subset S_i$ are quasi-compact
opens and $S = f_i^{-1}(V_{1, i}) \cup \ldots \cup f_i^{-1}(V_{n, i})$
then $S_{i'} = f_{i'i}^{-1}(V_{1, i}) \cup \ldots \cup f_{i'i}^{-1}(V_{n, i})$
for some $i' \geq i$.
\end{enumerate}
\end{lemma}

\begin{proof}
Choose $i_0 \in I$. Note that $I$ is nonempty as the limit is directed.
For convenience we write $S_0 = S_{i_0}$ and $i_0 = 0$.
Choose an affine open covering $S_0 = U_{1, 0} \cup \ldots \cup U_{m, 0}$.
Denote $U_{j, i} \subset S_i$ the inverse image of $U_{j, 0}$
under the transition morphism for $i \geq 0$.
Denote $U_j$ the inverse image of $U_{j, 0}$ in $S$.
Note that $U_j = \lim_i U_{j, i}$ is a limit of affine
schemes.

\medskip\noindent
We first prove the uniqueness statement: Let
$V_i \subset S_i$ and $V_{i'} \subset S_{i'}$
quasi-compact opens such that $f_i^{-1}(V_i) = f_{i'}^{-1}(V_{i'})$.
It suffices to show that $f_{i''i}^{-1}(V_i \cap U_{j, i''})$ and
$f_{i''i'}^{-1}(V_{i'} \cap U_{j, i''})$ become equal
for $i''$ large enough. Hence we reduce to the case
of a limit of affine schemes. In this case write
$S = \Spec(R)$ and $S_i = \Spec(R_i)$ for all $i \in I$.
We may write $V_i = S_i \setminus V(h_1, \ldots, h_m)$
and $V_{i'} = S_{i'} \setminus V(g_1, \ldots, g_n)$.
The assumption means that the ideals
$\sum g_jR$ and $\sum h_jR$ have the same radical
in $R$. This means that $g_j^N = \sum a_{jj'}h_{j'}$ and
$h_j^N = \sum b_{jj'} g_{j'}$ for some $N \gg 0$ and $a_{jj'}$
and $b_{jj'}$ in $R$.
Since $R = \colim_i R_i$ we can chose an index
$i'' \geq i$ such that the equations
$g_j^N = \sum a_{jj'}h_{j'}$ and
$h_j^N = \sum b_{jj'} g_{j'}$ hold in $R_{i''}$ for some
$a_{jj'}$ and $b_{jj'}$ in $R_{i''}$. This implies that
the ideals $\sum g_jR_{i''}$ and $\sum h_jR_{i''}$ have the same radical
in $R_{i''}$ as desired.

\medskip\noindent
We prove existence: If $S_0$ is affine, then $S_i = \Spec(R_i)$ for all
$i \geq 0$ and $S = \Spec(R)$ with $R = \colim R_i$. Then
$V = S \setminus V(g_1, \ldots, g_n)$ for some $g_1, \ldots, g_n \in R$.
Choose any $i$ large enough so that each of the $g_j$ comes from an
element $g_{j, i} \in R_i$ and take
$V_i = S_i \setminus V(g_{1, i}, \ldots, g_{n, i})$.
If $S_0$ is general, then the opens $V \cap U_j$
are quasi-compact because $S$ is quasi-separated. Hence by the
affine case we see that for each $j = 1, \ldots, m$
there exists an $i_j \in I$ and a quasi-compact open
$V_{i_j} \subset U_{j, i_j}$ whose inverse image in $U_j$
is $V \cap U_j$. Set $i = \max(i_1, \ldots, i_m)$
and let $V_i = \bigcup f_{ii_j}^{-1}(V_{i_j})$.

\medskip\noindent
The statement on coverings follows from the uniqueness statement
for the opens $V_{1, i} \cup \ldots \cup V_{n, i}$ and $S_i$ of $S_i$.
\end{proof}


\begin{lemma}
\label{limits-lemma-limit-quasi-affine}
In Situation \ref{limits-situation-descent} if $S$ is quasi-affine, then
for some $i_0 \in I$ the schemes $S_i$ for $i \geq i_0$ are quasi-affine.
\end{lemma}

\begin{proof}
Choose $i_0 \in I$. Note that $I$ is nonempty as the limit is directed.
For convenience we write $S_0 = S_{i_0}$ and $i_0 = 0$.
Let $s \in S$. We may choose an affine open
$U_0 \subset S_0$ containing $f_0(s)$. Since $S$ is quasi-affine
we may choose an element $a \in \Gamma(S, \mathcal{O}_S)$ such
that $s \in D(a) \subset f_0^{-1}(U_0)$, and such that
$D(a)$ is affine. By Lemma \ref{limits-lemma-descend-section}
there exists an $i \geq 0$ such that $a$
comes from an element $a_i \in \Gamma(S_i, \mathcal{O}_{S_i})$.
For any index $j \geq i$ we denote $a_j$
the image of $a_i$ in the global sections of the
structure sheaf of $S_j$.
Consider the opens $D(a_j) \subset S_j$
and $U_j = f_{j0}^{-1}(U_0)$. Note that
$U_j$ is affine and $D(a_j)$ is a quasi-compact open of $S_j$,
see Properties, Lemma \href{https://stacks.math.columbia.edu/tag/01PV}{lemma affine cap s open (Tag 01PV)}
for example. Hence we may apply Lemma \ref{limits-lemma-descend-opens} to the opens
$U_j$ and $U_j \cup D(a_j)$ to conclude that
$D(a_j) \subset U_j$ for some  $j \geq i$.
For such an index $j$ we see that $D(a_j) \subset S_j$ is an affine open
(because $D(a_j)$ is a standard affine open of the affine open $U_j$)
containing the image $f_j(s)$.

\medskip\noindent
We conclude that for every $s \in S$ there exist
an index $i \in I$, and a global section
$a \in \Gamma(S_i, \mathcal{O}_{S_i})$
such that $D(a) \subset S_i$ is an affine open
containing $f_i(s)$. Because $S$ is quasi-compact we
may choose a single index $i \in I$ and global sections
$a_1, \ldots, a_m \in \Gamma(S_i, \mathcal{O}_{S_i})$
such that each $D(a_j) \subset S_i$ is affine open
and such that $f_i : S \to S_i$ has image contained
in the union $W_i = \bigcup_{j = 1, \ldots, m} D(a_j)$.
For $i' \geq i$ set $W_{i'} = f_{i'i}^{-1}(W_i)$.
Since $f_i^{-1}(W_i)$ is all of $S$ we see
(by Lemma \ref{limits-lemma-descend-opens} again)
that for a suitable $i' \geq i$ we
have $S_{i'} = W_{i'}$. Thus we may replace $i$ by
$i'$ and assume that $S_i = \bigcup_{j = 1, \ldots, m} D(a_j)$.
This implies that $\mathcal{O}_{S_i}$ is an ample invertible
sheaf on $S_i$ (see Properties, Definition \ref{properties-definition-ample})
and hence that $S_i$ is quasi-affine, see
Properties, Lemma \href{https://stacks.math.columbia.edu/tag/01QE}{lemma quasi affine O ample (Tag 01QE)}.
Hence we win.
\end{proof}


\begin{lemma}
\label{limits-lemma-descend-section}
In Situation \ref{limits-situation-descent}.
Suppose that $\mathcal{F}_0$ is a quasi-coherent sheaf on $S_0$.
Set $\mathcal{F}_i = f_{i0}^*\mathcal{F}_0$ for $i \geq 0$ and set
$\mathcal{F} = f_0^*\mathcal{F}_0$.
Then
$$
\Gamma(S, \mathcal{F}) = \colim_{i \geq 0} \Gamma(S_i, \mathcal{F}_i)
$$
\end{lemma}

\begin{proof}
Write $\mathcal{A}_j = f_{i0, *} \mathcal{O}_{S_i}$.
This is a quasi-coherent sheaf of $\mathcal{O}_{S_0}$-algebras
(see Morphisms, Lemma \href{https://stacks.math.columbia.edu/tag/01SA}{lemma affine equivalence algebras (Tag 01SA)})
and $S_i$ is the relative spectrum of $\mathcal{A}_i$ over $S_0$.
In the proof of Lemma \ref{limits-lemma-directed-inverse-system-has-limit}
we constructed $S$ as the relative spectrum of
$\mathcal{A} = \colim_{i \geq 0} \mathcal{A}_i$
over $S_0$. Set
$$
\mathcal{M}_i = \mathcal{F}_0 \otimes_{\mathcal{O}_{S_0}} \mathcal{A}_i
$$
and
$$
\mathcal{M} = \mathcal{F}_0 \otimes_{\mathcal{O}_{S_0}} \mathcal{A}.
$$
Then we have $f_{i0, *} \mathcal{F}_i = \mathcal{M}_i$
and $f_{0, *}\mathcal{F} = \mathcal{M}$. Since $\mathcal{A}$
is the colimit of the sheaves $\mathcal{A}_i$ and since tensor
product commutes with directed colimits, we conclude that
$\mathcal{M} = \colim_{i \geq 0} \mathcal{M}_i$.
Since $S_0$ is quasi-compact and quasi-separated we see that
\begin{eqnarray*}
\Gamma(S, \mathcal{F})
& = &
\Gamma(S_0, \mathcal{M}) \\
& = &
\Gamma(S_0, \colim_{i \geq 0} \mathcal{M}_i) \\
& = &
\colim_{i \geq 0} \Gamma(S_0, \mathcal{M}_i) \\
& = &
\colim_{i \geq 0} \Gamma(S_i, \mathcal{F}_i)
\end{eqnarray*}
see Sheaves, Lemma \href{https://stacks.math.columbia.edu/tag/009F}{lemma directed colimits sections (Tag 009F)} and
Topology, Lemma \href{https://stacks.math.columbia.edu/tag/0069}{lemma topology quasi separated scheme (Tag 0069)}
for the middle equality.
\end{proof}


\begin{lemma}
\label{limits-lemma-quasi-affine-finite-type-over-Z}
Let $W$ be a quasi-affine scheme of finite type over
$\mathbf{Z}$. Suppose $W \to \Spec(R)$ is an
open immersion into an affine scheme. There exists a
finite type $\mathbf{Z}$-algebra $A \subset R$
which induces an open immersion $W \to \Spec(A)$.
Moreover, $R$ is the directed colimit of such subalgebras.
\end{lemma}

\begin{proof}
Choose an affine open covering $W = \bigcup_{i = 1, \ldots, n} W_i$
such that each $W_i$ is a standard affine open in $\Spec(R)$.
In other words, if we write $W_i = \Spec(R_i)$
then $R_i = R_{f_i}$ for some $f_i \in R$.
Choose finitely many $x_{ij} \in R_i$ which generate
$R_i$ over $\mathbf{Z}$.
Pick an $N \gg 0$ such that each $f_i^Nx_{ij}$ comes from an
element of $R$, say $y_{ij} \in R$.
Set $A$ equal to the $\mathbf{Z}$-algebra generated by
the $f_i$ and the $y_{ij}$ and (optionally) finitely many
additional elements of $R$. Then $A$ works. Details omitted.
\end{proof}


\begin{lemma}
\label{limits-lemma-diagram}
Suppose given a cartesian diagram of rings
$$
\xymatrix{
B \ar[r]_s & R \\
B'\ar[u] \ar[r] & R' \ar[u]_t
}
$$
Let $W' \subset \Spec(R')$ be an open of
the form $W' = D(f_1) \cup \ldots \cup D(f_n)$
such that $t(f_i) = s(g_i)$ for some $g_i \in B$
and $B_{g_i} \cong R_{s(g_i)}$. Then $B' \to R'$
induces an open immersion of $W'$ into $\Spec(B')$.
\end{lemma}

\begin{proof}
Set $h_i = (g_i, f_i) \in B'$. More on Algebra,
Lemma \ref{more-algebra-lemma-diagram-localize} shows that
$(B')_{h_i} \cong (R')_{f_i}$ as desired.
\end{proof}


\begin{lemma}
\label{more-algebra-lemma-diagram-localize}
Suppose given a cartesian diagram of rings
$$
\xymatrix{
R &
R' \ar[l]^t \\
B \ar[u]_s &
B'\ar[u] \ar[l]
}
$$
i.e., $B' = B \times_R R'$. If $h \in B'$ corresponds to $g \in B$
and $f \in R'$ such that $s(g) = t(f)$, then the diagram
$$
\xymatrix{
R_{s(g)} = R_{t(f)} &
(R')_f \ar[l]^-t \\
B_g \ar[u]_s &
(B')_h \ar[u] \ar[l]
}
$$
is cartesian too.
\end{lemma}

\begin{proof}
The equality $B' = B \times_R R'$ tells us that
$$
0 \to B' \to B \oplus R' \xrightarrow{s, -t} R
$$
is an exact sequence of $B'$-modules. We have $B_g = B_h$,
$R'_f = R'_h$, and $R_{s(g)} = R_{t(f)} = R_h$ as $B'$-modules.
By exactness of localization
(Algebra, Proposition \href{https://stacks.math.columbia.edu/tag/00CS}{proposition localization exact (Tag 00CS)})
we find that
$$
0 \to B'_h \to B_g \oplus R'_f \xrightarrow{s, -t} R_{s(g)} = R_{t(f)}
$$
is an exact sequence. This proves the lemma.
\end{proof}


\begin{lemma}
\label{limits-lemma-directed-inverse-system-has-limit}
Let $I$ be a directed set. Let $(S_i, f_{ii'})$ be an
inverse system of schemes over $I$. If all the morphisms
$f_{ii'} : S_i \to S_{i'}$ are affine, then the limit $S = \lim_i S_i$ exists
in the category of schemes. Moreover,
\begin{enumerate}
\item each of the morphisms $f_i : S \to S_i$ is affine,
\item for an element $0 \in I$ and any open subscheme $U_0 \subset S_0$
we have
$$
f_0^{-1}(U_0) = \lim_{i \geq 0} f_{i0}^{-1}(U_0)
$$
in the category of schemes.
\end{enumerate}
\end{lemma}

\begin{proof}
Choose an element $0 \in I$. Note that $I$ is nonempty as the limit is
directed. For every $i \geq 0$ consider the quasi-coherent sheaf of
$\mathcal{O}_{S_0}$-algebras $\mathcal{A}_i = f_{i0, *}\mathcal{O}_{S_i}$.
Recall that $S_i = \underline{\Spec}_{S_0}(\mathcal{A}_i)$,
see Morphisms, Lemma \href{https://stacks.math.columbia.edu/tag/01S8}{lemma characterize affine (Tag 01S8)}.
Set $\mathcal{A} = \colim_{i \geq 0} \mathcal{A}_i$.
This is a quasi-coherent sheaf of $\mathcal{O}_{S_0}$-algebras,
see Schemes, Section \href{https://stacks.math.columbia.edu/tag/01LA}{section quasi coherent (Tag 01LA)}.
Set $S = \underline{\Spec}_{S_0}(\mathcal{A})$.
By Morphisms, Lemma \href{https://stacks.math.columbia.edu/tag/01SA}{lemma affine equivalence algebras (Tag 01SA)}
we get for $i \geq 0$ morphisms $f_i : S \to S_i$ compatible with
the transition morphisms. Note that the morphisms $f_i$ are
affine by Morphisms, Lemma \href{https://stacks.math.columbia.edu/tag/01SG}{lemma affine permanence (Tag 01SG)} for example.
By Lemma \href{https://stacks.math.columbia.edu/tag/01YW}{lemma directed inverse system affine schemes has limit (Tag 01YW)} above
we see that for any affine open $U_0 \subset S_0$ the
inverse image $U = f_0^{-1}(U_0) \subset S$ is the limit of the
system of opens $U_i = f_{i0}^{-1}(U_0)$, $i \geq 0$ in the
category of schemes.

\medskip\noindent
Let $T$ be a scheme. Let $g_i : T \to S_i$ be a compatible system
of morphisms. To show that $S = \lim_i S_i$ we have
to prove there is a unique morphism $g : T \to S$ with
$g_i = f_i \circ g$ for all $i \in I$.
For every $t \in T$ there exists an affine open
$U_0 \subset S_0$ containing $g_0(t)$. Let $V \subset g_0^{-1}(U_0)$
be an affine open neighbourhood containing $t$.
By the remarks above we obtain a unique morphism
$g_V : V \to U = f_0^{-1}(U_0)$ such that $f_i \circ g_V = g_i|_{U_i}$
for all $i$. The open sets $V \subset T$ so constructed form
a basis for the topology of $T$. The morphisms $g_V$ glue to a morphism
$g : T \to S$ because of the uniqueness property. This gives the
desired morphism $g : T \to S$.

\medskip\noindent
The final statement is clear from the construction of the limit above.
\end{proof}


\begin{lemma}
\label{limits-lemma-descend-finite-presentation}
Let $I$ be a directed set.
Let $(S_i, f_{ii'})$ be an inverse system of schemes over $I$.
Assume
\begin{enumerate}
\item the morphisms $f_{ii'} : S_i \to S_{i'}$ are affine,
\item the schemes $S_i$ are quasi-compact and quasi-separated.
\end{enumerate}
Let $S = \lim_i S_i$. Then we have the following:
\begin{enumerate}
\item For any morphism of finite presentation $X \to S$
there exists an index $i \in I$ and a morphism of finite
presentation $X_i \to S_i$ such that $X \cong X_{i, S}$ as
schemes over $S$.
\item Given an index $i \in I$, schemes
$X_i$, $Y_i$ of finite presentation over $S_i$, and a morphism
$\varphi : X_{i, S} \to Y_{i, S}$ over $S$, there exists an index
$i' \geq i$ and a morphism
$\varphi_{i'} : X_{i, S_{i'}} \to Y_{i, S_{i'}}$
whose base change to $S$ is $\varphi$.
\item Given an index $i \in I$, schemes $X_i$, $Y_i$ of finite presentation
over $S_i$ and a pair of morphisms $\varphi_i, \psi_i : X_i \to Y_i$
whose base changes $\varphi_{i, S} = \psi_{i, S}$ are equal,
there exists an index $i' \geq i$ such that
$\varphi_{i, S_{i'}} = \psi_{i, S_{i'}}$.
\end{enumerate}
In other words, the category of schemes of finite presentation over
$S$ is the colimit over $I$ of the categories of schemes of finite
presentation over $S_i$.
\end{lemma}

\begin{proof}
In case each of the schemes $S_i$ is affine, and we consider
only affine schemes of finite presentation over $S_i$, resp.\ $S$
this lemma is equivalent to
Algebra, Lemma \ref{algebra-lemma-colimit-category-fp-algebras}.
We claim that the affine case implies the lemma in general.

\medskip\noindent
Let us prove (3). Suppose given an index $i \in I$, schemes
$X_i$, $Y_i$ of finite presentation over $S_i$ and a pair of morphisms
$\varphi_i, \psi_i : X_i \to Y_i$. Assume that the base changes are
equal: $\varphi_{i, S} = \psi_{i, S}$. We will use the notation
$X_{i'} = X_{i, S_{i'}}$ and $Y_{i'} = Y_{i, S_{i'}}$ for
$i' \geq i$. We also set $X = X_{i, S}$ and $Y = Y_{i, S}$.
Note that according to Lemma \ref{limits-lemma-scheme-over-limit} we have
$X = \lim_{i' \geq i} X_{i'}$ and similarly for $Y$.
Additionally we denote $\varphi_{i'}$ and $\psi_{i'}$
(resp.\ $\varphi$ and $\psi$)
the base change of $\varphi_i$ and $\psi_i$ to $S_{i'}$
(resp.\ $S$). So our assumption means that $\varphi = \psi$.
Since $Y_i$ and $X_i$ are of finite presentation
over $S_i$, and since $S_i$ is quasi-compact and quasi-separated, also
$X_i$ and $Y_i$ are quasi-compact and quasi-separated
(see Morphisms,
Lemma \href{https://stacks.math.columbia.edu/tag/01TY}{lemma finite presentation quasi compact quasi separated (Tag 01TY)}).
Hence we may choose a finite affine open covering
$Y_i = \bigcup V_{j, i}$ such that each $V_{j, i}$ maps into
an affine open of $S_i$. As above, denote $V_{j, i'}$ the inverse
image of $V_{j, i}$ in $Y_{i'}$ and $V_j$ the inverse image in $Y$.
The immersions $V_{j, i'} \to Y_{i'}$ are quasi-compact, and the inverse images
$U_{j, i'} = \varphi_i^{-1}(V_{j, i'})$ and
$U_{j, i'}' = \psi_i^{-1}(V_{j, i'})$
are quasi-compact opens of $X_{i'}$. By assumption the inverse images of
$V_j$ under $\varphi$ and $\psi$ in $X$ are equal.
Hence by Lemma \ref{limits-lemma-descend-opens}
there exists an index $i' \geq i$ such that
of $U_{j, i'} = U_{j, i'}'$ in $X_{i'}$.
Choose an finite affine open covering
$U_{j, i'} = U_{j, i'}' = \bigcup W_{j, k, i'}$
which induce coverings $U_{j, i''} = U_{j, i''}' = \bigcup W_{j, k, i''}$
for all $i'' \geq i'$.
By the affine case there exists
an index $i''$ such that
$\varphi_{i''}|_{W_{j, k, i''}} = \psi_{i''}|_{W_{j, k, i''}}$
for all $j, k$. Then $i''$ is an index such that
$\varphi_{i''} = \psi_{i''}$ and (3) is proved.

\medskip\noindent
Let us prove (2). Suppose given an index $i \in I$, schemes
$X_i$, $Y_i$ of finite presentation over $S_i$ and a morphism
$\varphi : X_{i, S} \to Y_{i, S}$. We will use the notation
$X_{i'} = X_{i, S_{i'}}$ and $Y_{i'} = Y_{i, S_{i'}}$ for
$i' \geq i$. We also set $X = X_{i, S}$ and $Y = Y_{i, S}$.
Note that according to Lemma \ref{limits-lemma-scheme-over-limit} we have
$X = \lim_{i' \geq i} X_{i'}$ and similarly for $Y$.
Since $Y_i$ and $X_i$ are of finite presentation
over $S_i$, and since $S_i$ is quasi-compact and quasi-separated, also
$X_i$ and $Y_i$ are quasi-compact and quasi-separated
(see Morphisms,
Lemma \href{https://stacks.math.columbia.edu/tag/01TY}{lemma finite presentation quasi compact quasi separated (Tag 01TY)}).
Hence we may choose a finite affine open covering
$Y_i = \bigcup V_{j, i}$ such that each $V_{j, i}$ maps into
an affine open of $S_i$. As above, denote $V_{j, i'}$ the inverse
image of $V_{j, i}$ in $Y_{i'}$ and $V_j$ the inverse image in $Y$.
The immersions $V_j \to Y$ are quasi-compact, and the inverse images
$U_j = \varphi^{-1}(V_j)$ are quasi-compact opens of $X$.
Hence by Lemma \ref{limits-lemma-descend-opens} there exists an index
$i' \geq i$ and quasi-compact opens $U_{j, i'}$ of $X_{i'}$
whose inverse image in $X$ is $U_j$. Choose an finite affine open covering
$U_{j, i'} = \bigcup W_{j, k, i'}$ which induce affine open
coverings $U_{j, i''} = \bigcup W_{j, k, i''}$
for all $i'' \geq i'$ and an affine open covering
$U_j = \bigcup W_{j, k}$. By the affine case there exists
an index $i''$ and morphisms
$\varphi_{j, k, i''} : W_{j, k, i''} \to V_{j, i''}$
such that
$\varphi|_{W_{j, k}} = \varphi_{j, k, i'', S}$ for all $j, k$.
By part (3) proved above, there is a further index $i''' \geq i''$
such that
$$
\varphi_{j_1, k_1, i'', S_{i'''}}|_{W_{j_1, k_1, i'''} \cap W_{j_2, k_2, i'''}}
=
\varphi_{j_2, k_2, i'', S_{i'''}}|_{W_{j_1, k_1, i'''} \cap W_{j_2, k_2, i'''}}
$$
for all $j_1, j_2, k_1, k_2$. Then $i'''$ is an index such that
there exists a morphism $\varphi_{i'''} : X_{i'''} \to Y_{i'''}$
whose base change to $S$ gives $\varphi$. Hence (2) holds.

\medskip\noindent
Let us prove (1). Suppose given a scheme $X$ of finite presentation
over $S$. Since $X$ is of finite presentation
over $S$, and since $S$ is quasi-compact and quasi-separated, also
$X$ is quasi-compact and quasi-separated
(see Morphisms,
Lemma \href{https://stacks.math.columbia.edu/tag/01TY}{lemma finite presentation quasi compact quasi separated (Tag 01TY)}).
Choose a finite affine open covering $X = \bigcup U_j$
such that each $U_j$ maps into an affine open $V_j \subset S$.
Denote $U_{j_1j_2} = U_{j_1} \cap U_{j_2}$ and
$U_{j_1j_2j_3} = U_{j_1} \cap U_{j_2} \cap U_{j_3}$.
By Lemmas \ref{limits-lemma-descend-opens} and \ref{limits-lemma-limit-affine}
we can find an index $i_1$ and affine opens $V_{j, i_1} \subset S_{i_1}$
such that each $V_j$ is the inverse of this in $S$.
Let $V_{j, i}$ be the inverse image of $V_{j, i_1}$ in $S_i$ for
$i \geq i_1$. By the affine case we may find an index $i_2 \geq i_1$ and
affine schemes $U_{j, i_2} \to V_{j, i_2}$ such
that $U_j = S \times_{S_{i_2}} U_{j, i_2}$ is the base change.
Denote $U_{j, i} = S_i \times_{S_{i_2}} U_{j, i_2}$ for $i \geq i_2$.
By Lemma \ref{limits-lemma-descend-opens} there exists an index
$i_3 \geq i_2$ and open subschemes
$W_{j_1, j_2, i_3} \subset U_{j_1, i_3}$
whose base change to $S$ is equal to $U_{j_1j_2}$.
Denote $W_{j_1, j_2, i} = S_i \times_{S_{i_3}} W_{j_1, j_2, i_3}$
for $i \geq i_3$. By part (2) shown above there exists an index
$i_4 \geq i_3$ and morphisms
$\varphi_{j_1, j_2, i_4} : W_{j_1, j_2, i_4} \to W_{j_2, j_1, i_4}$
whose base change to $S$ gives the identity morphism
$U_{j_1j_2} = U_{j_2j_1}$ for all $j_1, j_2$.
For all $i \geq i_4$ denote
$\varphi_{j_1, j_2, i} = \text{id}_S \times \varphi_{j_1, j_2, i_4}$
the base change. We claim that for some $i_5 \geq i_4$ the system
$((U_{j, i_5})_j, (W_{j_1, j_2, i_5})_{j_1, j_2},
(\varphi_{j_1, j_2, i_5})_{j_1, j_2})$ forms a glueing datum
as in Schemes, Section \href{https://stacks.math.columbia.edu/tag/01JA}{section glueing schemes (Tag 01JA)}.
In order to see this we have to verify that for $i$ large enough
we have
$$
\varphi_{j_1, j_2, i}^{-1}(W_{j_2, j_1, i} \cap W_{j_2, j_3, i})
=
W_{j_1, j_2, i} \cap W_{j_1, j_3, i}
$$
and that for large enough $i$ the cocycle condition holds.
The first condition follows from Lemma \ref{limits-lemma-descend-opens}
and the fact that $U_{j_2j_1j_3} = U_{j_1j_2j_3}$.
The second from part (3) of the lemma proved above and the fact
that the cocycle condition holds for the maps
$\text{id} : U_{j_1j_2} \to U_{j_2j_1}$.
Ok, so now we can use Schemes, Lemma \href{https://stacks.math.columbia.edu/tag/01JC}{lemma glue schemes (Tag 01JC)}
to glue the system
$((U_{j, i_5})_j, (W_{j_1, j_2, i_5})_{j_1, j_2},
(\varphi_{j_1, j_2, i_5})_{j_1, j_2})$ to get a scheme
$X_{i_5} \to S_{i_5}$. By construction the base change of
$X_{i_5}$ to $S$ is formed by glueing the open affines
$U_j$ along the opens $U_{j_1} \leftarrow U_{j_1j_2} \rightarrow U_{j_2}$.
Hence $S \times_{S_{i_5}} X_{i_5} \cong X$ as desired.
\end{proof}


\begin{lemma}
\label{limits-lemma-limit-affine}
In Situation \ref{limits-situation-descent} if $S$ is affine,
then for some $i_0 \in I$ the schemes $S_i$ for $i \geq i_0$
are affine.
\end{lemma}

\begin{proof}
By Lemma \ref{limits-lemma-limit-quasi-affine} we may assume that $S_0$ is
quasi-affine for some $0 \in I$. Set $R_0 = \Gamma(S_0, \mathcal{O}_{S_0})$.
Then $S_0$ is a quasi-compact open of $T_0 = \Spec(R_0)$. Denote
$j_0 : S_0 \to T_0$ the corresponding quasi-compact open immersion.
For $i \geq 0$ set $\mathcal{A}_i = f_{i0, *}\mathcal{O}_{S_i}$.
Since $f_{i0}$ is affine we see that
$S_i = \underline{\Spec}_{S_0}(\mathcal{A}_i)$.
Set $T_i = \underline{\Spec}_{T_0}(j_{0, *}\mathcal{A}_i)$.
Then $T_i \to T_0$ is affine, hence $T_i$ is affine. Thus
$T_i$ is the spectrum of
$$
R_i = \Gamma(T_0, j_{0, *}\mathcal{A}_i) = \Gamma(S_0, \mathcal{A}_i) =
\Gamma(S_i, \mathcal{O}_{S_i}).
$$
Write $S = \Spec(R)$. We have $R = \colim_i R_i$
by Lemma \ref{limits-lemma-descend-section}.
Hence also $S = \lim_i T_i$. As formation of the relative spectrum commutes
with base change, the inverse image
of the open $S_0 \subset T_0$ in $T_i$ is $S_i$.
Let $Z_0 = T_0 \setminus S_0$ and let $Z_i \subset T_i$
be the inverse image of $Z_0$. As $S_i = T_i \setminus Z_i$, it suffices
to show that $Z_i$ is empty for some $i$. Assume $Z_i$ is nonempty for all
$i$ to get a contradiction. By Lemma \ref{limits-lemma-limit-closed-nonempty}
there exists a point $s$ of $S = \lim T_i$ which maps to a point of $Z_i$
for every $i$. But $S = \lim_i S_i$, and hence we arrive at a contradiction
by Lemma \ref{limits-lemma-topology-limit}.
\end{proof}


\begin{lemma}
\label{limits-lemma-limit-closed-nonempty}
In Situation \ref{limits-situation-descent}.
Suppose for each $i$ we are given a nonempty closed subset
$Z_i \subset S_i$ with $f_{i'i}(Z_{i'}) \subset Z_i$ for all
$i' \geq i$.
Then there exists a point $s \in S$ with $f_i(s) \in Z_i$ for
all $i$.
\end{lemma}

\begin{proof}
Let $Z_i \subset S_i$ also denote the reduced closed subscheme
associated to $Z_i$, see Schemes,
Definition \ref{schemes-definition-reduced-induced-scheme}.
A closed immersion is affine, and a composition of affine
morphisms is affine (see
Morphisms, Lemmas \href{https://stacks.math.columbia.edu/tag/01SE}{lemma closed immersion affine (Tag 01SE)}
and \href{https://stacks.math.columbia.edu/tag/01SC}{lemma composition affine (Tag 01SC)}), and hence $Z_{i'} \to S_i$ is
affine when $i' \geq i$. We conclude that the morphism
$f_{i'i} : Z_{i'} \to Z_i$ is affine by
Morphisms, Lemma \href{https://stacks.math.columbia.edu/tag/01SG}{lemma affine permanence (Tag 01SG)}.
Each of the schemes $Z_i$ is quasi-compact as a closed
subscheme of a quasi-compact scheme. Hence we may apply
Lemma \ref{limits-lemma-limit-nonempty} to see that
$Z = \lim_i Z_i$ is nonempty. Since there is a
canonical morphism $Z \to S$ we win.
\end{proof}


\begin{lemma}
\label{limits-lemma-limit-nonempty}
Let $S = \lim S_i$ be the limit of a directed inverse system
of schemes with affine transition morphisms
(Lemma \ref{limits-lemma-directed-inverse-system-has-limit}).
If all the schemes $S_i$ are nonempty and quasi-compact,
then the limit $S = \lim_i S_i$ is nonempty.
\end{lemma}

\begin{proof}
Choose $0 \in I$. Note that $I$ is nonempty as the limit is directed.
Choose an affine open covering $S_0 = \bigcup_{j = 1, \ldots, m} U_j$.
Since $I$ is directed there exists a $j \in \{1, \ldots, m\}$
such that $f_{i0}^{-1}(U_j) \not = \emptyset$ for all
$i \geq 0$. Hence $\lim_{i \geq 0} f_{i0}^{-1}(U_j)$ is not
empty since a directed colimit of nonzero rings is nonzero
(because $1 \not = 0$). As $\lim_{i \geq 0} f_{i0}^{-1}(U_j)$
is an open subscheme of the limit we win.
\end{proof}


\begin{lemma}
\label{limits-lemma-topology-limit}
In Situation \ref{limits-situation-descent}.
\begin{enumerate}
\item We have $S_{set} = \lim_i S_{i, set}$ where $S_{set}$
indicates the underlying set of the scheme $S$.
\item We have $S_{top} = \lim_i S_{i, top}$ where $S_{top}$
indicates the underlying topological space of the scheme $S$.
\item If $s, s' \in S$ and $s'$ is not a specialization of $s$
then for some $i \in I$ the image $s'_i \in S_i$ of $s'$ is not
a specialization of the image $s_i \in S_i$ of $s$.
\item Add more easy facts on topology of $S$ here.
(Requirement: whatever is added should be easy in the affine case.)
\end{enumerate}
\end{lemma}

\begin{proof}
Part (1) is a special case of Lemma \href{https://stacks.math.columbia.edu/tag/0CUE}{lemma inverse limit sets (Tag 0CUE)}.

\medskip\noindent
Part (2) is a special case of Lemma \href{https://stacks.math.columbia.edu/tag/0CUF}{lemma inverse limit top (Tag 0CUF)}.

\medskip\noindent
Part (3) is a special case of Lemma \href{https://stacks.math.columbia.edu/tag/0CUG}{lemma inverse limit irreducibles (Tag 0CUG)}.
\end{proof}


\begin{lemma}
\label{limits-lemma-scheme-over-limit}
Let $I$ be a directed set.
Let $(S_i, f_{ii'})$ be an inverse system of schemes over $I$.
Assume all the morphisms $f_{ii'} : S_i \to S_{i'}$ are affine,
Let $S = \lim_i S_i$. Let $0 \in I$.
Suppose that $T$ is a scheme over $S_0$.
Then
$$
T \times_{S_0} S = \lim_{i \geq 0} T \times_{S_0} S_i
$$
\end{lemma}

\begin{proof}
The right hand side is a scheme by
Lemma \ref{limits-lemma-directed-inverse-system-has-limit}.
The equality is formal, see
Categories, Lemma \href{https://stacks.math.columbia.edu/tag/002M}{lemma colimits commute (Tag 002M)}.
\end{proof}


\begin{lemma}
\label{algebra-lemma-colimit-category-fp-algebras}
Suppose that $R = \colim_{\lambda \in \Lambda} R_\lambda$ is a directed colimit
of rings. Then the category of finitely presented $R$-algebras is
the colimit of the categories of finitely presented $R_\lambda$-algebras.
More precisely
\begin{enumerate}
\item Given a finitely presented $R$-algebra $A$ there exists a
$\lambda \in \Lambda$ and a finitely presented $R_\lambda$-algebra
$A_\lambda$ such that $A \cong A_\lambda \otimes_{R_\lambda} R$.
\item Given a $\lambda \in \Lambda$, finitely presented
$R_\lambda$-algebras $A_\lambda, B_\lambda$, and an $R$-algebra map
$\varphi : A_\lambda \otimes_{R_\lambda} R \to B_\lambda \otimes_{R_\lambda} R$,
then there exists a $\mu \geq \lambda$ and an $R_\mu$-algebra map
$\varphi_\mu : A_\lambda \otimes_{R_\lambda} R_\mu \to
B_\lambda \otimes_{R_\lambda} R_\mu$
such that $\varphi = \varphi_\mu \otimes 1_R$.
\item Given a $\lambda \in \Lambda$, finitely presented
$R_\lambda$-algebras $A_\lambda, B_\lambda$, and $R_\lambda$-algebra maps
$\varphi_\lambda, \psi_\lambda : A_\lambda \to B_\lambda$
such that $\varphi \otimes 1_R = \psi \otimes 1_R$, then
$\varphi \otimes 1_{R_\mu} = \psi \otimes 1_{R_\mu}$ for some
$\mu \geq \lambda$.
\end{enumerate}
\end{lemma}

\begin{proof}
To prove (1) choose a presentation
$A = R[x_1, \ldots, x_n]/(f_1, \ldots, f_m)$.
We can choose a $\lambda \in \Lambda$ and elements
$f_{\lambda, j} \in R_\lambda[x_1, \ldots, x_n]$ mapping to
$f_j \in R[x_1, \ldots, x_n]$.
Then we simply let
$A_\lambda =
R_\lambda[x_1, \ldots, x_n]/(f_{\lambda, 1}, \ldots, f_{\lambda, m})$.

\medskip\noindent
Parts (2) and (3) follow from
Lemma \ref{algebra-lemma-algebra-map-property-in-colimit}.
\end{proof}


\begin{lemma}
\label{algebra-lemma-algebra-map-property-in-colimit}
Let $A$ be a ring and let $B, C$ be $A$-algebras.
Suppose that $R = \colim_{i \in I} R_i$ is a directed colimit
of $A$-algebras.
\begin{enumerate}
\item If $B$ is a finite type $A$-algebra, and $u, u' : B \to C$ are
$A$-algebra maps such that
$u \otimes 1 = u' \otimes 1 : B \otimes_A R \to C \otimes_A R$
then for some $i$ we have
$u \otimes 1 = u' \otimes 1 : B \otimes_A R_i \to C \otimes_A R_i$.
\item If $C$ is a finite type $A$-algebra and $u : B \to C$ is an
$A$-algebra map such that
$u \otimes 1 : B \otimes_A R \to C \otimes_A R$ is surjective, then
for some $i$ the map $u \otimes 1 : B \otimes_A R_i \to C \otimes_A R_i$
is surjective.
\item If $C$ is of finite presentation over $A$ and
$v : C \otimes_A R \to B \otimes_A R$ is an $R$-algebra map, then there
exists an $i$ and an $R_i$-algebra map
$v_i : C \otimes_A R_i \to B \otimes_A R_i$ such that
$v = v_i \otimes 1$.
\item If $B$ is a finite type $A$-algebra, $C$ is a finitely presented
$A$-algebra, and
$u \otimes 1 : B \otimes_A R \to C \otimes_A R$ is an isomorphism, then
for some $i$ the map $u \otimes 1 : B \otimes_A R_i \to C \otimes_A R_i$
is an isomorphism.
\end{enumerate}
\end{lemma}

\begin{proof}
To prove (1) assume $u$ is as in (1) and
let $x_1, \ldots, x_m \in B$ be generators. Since
$B \otimes_A R = \colim_i B \otimes_A R_i$
we may pick an $i \in I$ such that $u(x_j) \otimes 1 = u'(x_j) \otimes 1$
in $B \otimes_A R_i$, $j = 1, \ldots, m$.
For such an $i$ we have
$u \otimes 1 = u' \otimes 1 : B \otimes_A R_i \to C \otimes_A R_i$.

\medskip\noindent
To prove (2) assume $u \otimes 1$ surjective and
let $y_1, \ldots, y_m \in C$ be generators. Since
$B \otimes_A R = \colim_i B \otimes_A R_i$
we may pick an $i \in I$ and $z_j \in B \otimes_A R_i$, $j = 1, \ldots, m$
whose images in $C \otimes_A R$ equal $y_j \otimes 1$.
For such an $i$ the map $u \otimes 1 : B \otimes_A R_i \to C \otimes_A R_i$
is surjective.

\medskip\noindent
To prove (3) let $c_1, \ldots, c_m \in C$ be generators. Let
$K = \Ker(A[x_1, \ldots, x_m] \to N)$ where the map is given by
the rule $x_j \mapsto \sum c_j$. Let $f_1, \ldots, f_t$
be generators for $K$ as an ideal in $A[x_1, \ldots, x_m]$.
We think of $f_j = f_j(x_1, \ldots, x_m)$ as a polynomial.
Since $B \otimes_A R = \colim_i B \otimes_A R_i$
we may pick an $i \in I$ and $z_j \in B \otimes_A R_i$, $j = 1, \ldots, m$
whose images in $B \otimes_A R$ equal $v(c_j \otimes 1)$.
We want to use the $z_j$ to define a map
$v_i : C \otimes_A R_i \to B \otimes_A R_i$.
Since $K \otimes_A R_i \to R_i[x_1, \ldots, x_m] \to C \otimes_A R_i \to 0$
is a presentation, it suffices to check that
$\xi_s = f_j(z_1, \ldots, z_m)$ is
zero in $B \otimes_A R_i$ for each $s = 1, \ldots, t$. This may not
be the case, but since the image of $\xi_s$ in $B \otimes_A R$ is zero
we see that it will be the case after increasing $i$ a bit.

\medskip\noindent
To prove (4) assume $u \otimes 1$ is an isomorphism, that
$B$ is a finite type $A$-algebra, and that $C$ is a finitely presented
$A$-algebra. Let $v : B \otimes_A R \to C \otimes_A R$ be an inverse to
$u \otimes 1$. Let $v_i : C \otimes_A R_i \to B \otimes_A R_i$ be as
in part (3). Apply part (1) to see that, after increasing $i$ we have
$v_i \circ (u \otimes 1) = \text{id}_{B \otimes_R R_i}$ and
$(u \otimes 1) \circ v_i = \text{id}_{C \otimes_R R_i}$.
\end{proof}


\begin{proposition}
\label{proposition-approximate}
\begin{reference}
Our proof follows closely the proof given in \href{https://stacks.math.columbia.edu/bibliography}{Bibliography: CLO (Theorem 1.2.2)}.
\end{reference}
Let $X$ be a quasi-compact and quasi-separated algebraic space over
$\Spec(\mathbf{Z})$. There exist a directed set $I$
and an inverse system of algebraic spaces $(X_i, f_{ii'})$ over $I$
such that
\begin{enumerate}
\item the transition morphisms $f_{ii'}$ are affine
\item each $X_i$ is quasi-separated and of finite type over
$\mathbf{Z}$, and
\item $X = \lim X_i$.
\end{enumerate}
\end{proposition}

\begin{proof}
We apply Decent Spaces, Lemma
\ref{decent-spaces-lemma-filter-quasi-compact-quasi-separated}
to get open subspaces $U_p \subset X$, schemes $V_p$, and morphisms
$f_p : V_p \to U_p$ with properties as stated. Note that
$f_n : V_n \to U_n$ is an \'etale morphism of algebraic spaces
whose restriction to the inverse image of $T_n = (V_n)_{red}$ is an
isomorphism. Hence $f_n$ is an isomorphism, for example by
Morphisms of Spaces, Lemma
\href{https://stacks.math.columbia.edu/tag/05W5}{lemma etale universally injective open (Tag 05W5)}.
In particular $U_n$ is a quasi-compact and separated scheme.
Thus we can write $U_n = \lim U_{n, i}$ as a directed limit
of schemes of finite type over $\mathbf{Z}$ with affine transition
morphisms, see Limits, Proposition \href{https://stacks.math.columbia.edu/tag/01ZA}{proposition approximate (Tag 01ZA)}.
Thus, applying descending induction on $p$, we see that we have reduced
to the problem posed in the following paragraph.

\medskip\noindent
Here we have $U \subset X$, $U = \lim U_i$, $Z \subset X$, and
$f : V \to X$ with the following properties
\begin{enumerate}
\item $X$ is a quasi-compact and quasi-separated algebraic space,
\item $V$ is a quasi-compact and separated scheme,
\item $U \subset X$ is a quasi-compact open subspace,
\item $(U_i, g_{ii'})$ is a directed inverse system of
quasi-separated algebraic spaces
of finite type over $\mathbf{Z}$ with affine transition morphisms
whose limit is $U$,
\item $Z \subset X$ is a closed subspace such that $|X| = |U| \amalg |Z|$,
\item $f : V \to X$ is a surjective \'etale morphism such that
$f^{-1}(Z) \to Z$ is an isomorphism.
\end{enumerate}
Problem: Show that the conclusion of the proposition holds for $X$.

\medskip\noindent
Note that $W = f^{-1}(U) \subset V$ is a quasi-compact open subscheme
\'etale over $U$. Hence we may apply
Lemmas \ref{lemma-descend-finite-presentation} and \ref{lemma-descend-etale}
to find an index $0 \in I$ and an \'etale morphism $W_0 \to U_0$
of finite presentation whose base change to $U$ produces $W$. Setting
$W_i = W_0 \times_{U_0} U_i$ we see that $W = \lim_{i \geq 0} W_i$. After
increasing $0$ we may assume the $W_i$ are schemes, see
Lemma \ref{lemma-limit-is-scheme}.
Moreover, $W_i$ is of finite type over $\mathbf{Z}$.

\medskip\noindent
Apply Limits, Lemma \ref{limits-lemma-approximate} to
$W = \lim_{i \geq 0} W_i$ and the inclusion $W \subset V$. Replace $I$
by the directed set $J$ found in that lemma. This allows us
to write $V$ as a directed limit $V = \lim V_i$ of finite type schemes over
$\mathbf{Z}$ with affine transition maps such that each $V_i$ contains
$W_i$ as an open subscheme (compatible with transition morphisms).
For each $i$ we can form the push out
$$
\xymatrix{
W_i \ar[r] \ar[d]_\Delta & V_i \ar[d] \\
W_i \times_{U_i} W_i \ar[r] & R_i
}
$$
in the category of schemes. Namely, the left vertical and upper horizontal
arrows are open immersions of schemes. In other words, we can construct
$R_i$ as the glueing of $V_i$ and $W_i \times_{U_i} W_i$ along the common open
$W_i$ (see Schemes, Section \href{https://stacks.math.columbia.edu/tag/01JA}{section glueing schemes (Tag 01JA)}). Note that
the \'etale projection maps $W_i \times_{U_i} W_i \to W_i$ extend
to \'etale morphisms $s_i, t_i : R_i \to V_i$. It is clear that the
morphism $j_i = (t_i, s_i) : R_i \to V_i \times V_i$ is an \'etale
equivalence relation on $V_i$. Note that $W_i \times_{U_i} W_i$ is
quasi-compact (as $U_i$ is quasi-separated and $W_i$ quasi-compact)
and $V_i$ is quasi-compact, hence $R_i$ is quasi-compact. For
$i \geq i'$ the diagram
\begin{equation}
\label{equation-cartesian}
\vcenter{
\xymatrix{
R_i \ar[r] \ar[d]_{s_i} & R_{i'} \ar[d]^{s_{i'}} \\
V_i \ar[r] & V_{i'}
}
}
\end{equation}
is cartesian because
$$
(W_{i'} \times_{U_{i'}} W_{i'}) \times_{U_{i'}} U_i =
W_{i'} \times_{U_{i'}} U_i \times_{U_i} U_i \times_{U_{i'}} W_{i'} =
W_i \times_{U_i} W_i.
$$
Consider the algebraic space $X_i = V_i/R_i$ (see
Spaces, Theorem \ref{spaces-theorem-presentation}).
As $V_i$ is of finite type over $\mathbf{Z}$ and $R_i$ is quasi-compact
we see that $X_i$ is quasi-separated and of finite type over $\mathbf{Z}$
(see
Properties of Spaces, Lemma \href{https://stacks.math.columbia.edu/tag/07S4}{lemma quasi separated (Tag 07S4)}
and
Morphisms of Spaces, Lemmas
\href{https://stacks.math.columbia.edu/tag/040W}{lemma surjection from quasi compact (Tag 040W)} and
\href{https://stacks.math.columbia.edu/tag/040Y}{lemma finite type local (Tag 040Y)}).
As the construction of $R_i$ above is compatible
with transition morphisms, we obtain morphisms of algebraic spaces
$X_i \to X_{i'}$ for $i \geq i'$. The commutative diagrams
$$
\xymatrix{
V_i \ar[r] \ar[d] & V_{i'} \ar[d] \\
X_i \ar[r] & X_{i'}
}
$$
are cartesian as (\ref{equation-cartesian}) is cartesian, see
Groupoids, Lemma \ref{groupoids-lemma-criterion-fibre-product}.
Since $V_i \to V_{i'}$ is affine, this implies that $X_i \to X_{i'}$
is affine, see
Morphisms of Spaces, Lemma \href{https://stacks.math.columbia.edu/tag/03WG}{lemma affine local (Tag 03WG)}.
Thus we can form the limit $X' = \lim X_i$ by
Lemma \ref{lemma-directed-inverse-system-has-limit}.
We claim that $X \cong X'$ which finishes the proof of the proposition.

\medskip\noindent
Proof of the claim. Set $R = \lim R_i$.
By construction the algebraic space $X'$ comes
equipped with a surjective \'etale morphism $V \to X'$ such that
$$
V \times_{X'} V \cong R
$$
(use Lemma \ref{lemma-directed-inverse-system-has-limit}).
By construction $\lim W_i \times_{U_i} W_i = W \times_U W$ and $V = \lim V_i$
so that $R$ is the union of $W \times_U W$ and $V$ glued along $W$.
Property (6) implies the projections $V \times_X V \to V$ are isomorphisms
over $f^{-1}(Z) \subset V$. Hence the scheme $V \times_X V$ is the union
of the opens $\Delta_{V/X}(V)$ and $W \times_U W$ which intersect
along $\Delta_{W/X}(W)$. We conclude that there exists a unique isomorphism
$R \cong V \times_X V$ compatible with the projections to $V$.
Since $V \to X$ and $V \to X'$ are surjective \'etale we see that
$$
X = V/ V \times_X V = V/R = V/V \times_{X'} V = X'
$$
by Spaces, Lemma \href{https://stacks.math.columbia.edu/tag/0262}{lemma space presentation (Tag 0262)} and we win.
\end{proof}
\end{document}
